\documentclass{beamer}
\usetheme{Madrid}
\usecolortheme{default}
\usepackage[utf8]{inputenc}
\usepackage{amsmath}
\usepackage{amsfonts}
\usepackage{amssymb}
\usepackage{tikz}

\title{CRISC Domain D - 100 Hard Questions with Advanced Options}
\subtitle{Information Technology and Security - Expert Level Part 4}
\author{CRISC Mastery Series}
\date{\today}

\begin{document}
	
	\begin{frame}
		\titlepage
	\end{frame}
	
	\begin{frame}{Warning}
		\begin{alertblock}{Extreme Difficulty}
			Every question is scenario-based and asks for the BEST, FIRST, or MOST appropriate action.\\
			Every option is long and plausible; several are partly correct.\\
			Think like a risk practitioner advising the business, not a technician.
		\end{alertblock}
	\end{frame}
	
	\begin{frame}{Outline}
		\tableofcontents
	\end{frame}
	
	% ============================================================
	\section{1. Enterprise Architecture}
	% ============================================================
	
	\begin{frame}{Q301: Missing Gap Analysis}
		\begin{block}{Question}
			\small A risk practitioner reviewing an architecture program finds that the target-state design was approved without any gap analysis between the baseline and target architectures. What is the MOST significant risk?
		\end{block}
		\begin{enumerate}[(A)]
			\footnotesize
			\item The architecture board will be unable to meet again, because gap analysis is a mandatory precondition for scheduling any further review sessions
			\item The target architecture will be technically invalid, because baseline documentation is always more reliable than any future-state design
			\item Migration planning and funding will lack evidence of what must change, leading to underestimated cost, missed dependencies, and unmanaged transition risk
			\item Vendors will be prevented from bidding, because requirements cannot be expressed to the market without a completed gap register
		\end{enumerate}
	\end{frame}
	
	\begin{frame}{Q301: Explanation}
		\begin{block}{Correct Answer: C}
			\textbf{Why (C) is the best answer:}
			\begin{itemize}
				\item Gap analysis links baseline to target state
				\item It drives sequencing, cost, and dependency decisions
				\item Without it, transition risk stays hidden
			\end{itemize}
		\end{block}
		\begin{alertblock}{Critical Trap}
			A target state without gap analysis is an aspiration, not a plan.
		\end{alertblock}
	\end{frame}
	
	\begin{frame}{Q302: Conflicting Principles}
		\begin{block}{Question}
			\small Two approved principles conflict on a major initiative: ``buy before build'' and ``minimize exposure of regulated data to third parties.'' The sponsor wants a SaaS product that stores regulated customer data. What should the risk practitioner recommend FIRST?
		\end{block}
		\begin{enumerate}[(A)]
			\footnotesize
			\item Apply whichever principle has the lowest implementation cost, since principles are ranked by their financial impact on the project
			\item Escalate through architecture governance for an explicit, documented trade-off decision, with risk acceptance by the accountable business owner
			\item Reject the SaaS product, because a conflict between principles is always resolved in favor of the more restrictive principle
			\item Approve the product and add compensating controls after go-live, since principles are guidance and do not require formal governance review
		\end{enumerate}
	\end{frame}
	
	\begin{frame}{Q302: Explanation}
		\begin{block}{Correct Answer: B}
			\textbf{Why (B) is the best answer:}
			\begin{itemize}
				\item Conflicts need a governed decision, not a silent choice
				\item Trade-offs must be documented with rationale
				\item Risk acceptance belongs to the accountable owner
			\end{itemize}
		\end{block}
		\begin{alertblock}{Critical Trap}
			Principle conflicts are resolved by governance, not by convenience.
		\end{alertblock}
	\end{frame}
	
	\begin{frame}{Q303: Architecture Debt Prioritization}
		\begin{block}{Question}
			\small An architecture team reports that tactical integrations built to meet deadlines now generate most of the incident volume. Leadership asks whether to fund remediation. Which analysis BEST supports a risk-based decision?
		\end{block}
		\begin{enumerate}[(A)]
			\footnotesize
			\item Compare the remediation cost with the original project budget to determine whether the earlier shortcuts were cost-effective
			\item Rank integrations by age and retire the oldest first, since age is the strongest single predictor of failure and security weakness
			\item Defer remediation until incident volume exceeds the threshold written into the integration vendor's support contract
			\item Quantify the ongoing operational and security exposure of the debt against the cost and risk of remediation, then prioritize by business criticality
		\end{enumerate}
	\end{frame}
	
	\begin{frame}{Q303: Explanation}
		\begin{block}{Correct Answer: D}
			\textbf{Why (D) is the best answer:}
			\begin{itemize}
				\item Debt is a risk exposure, not just a cost
				\item Compare exposure against remediation effort
				\item Prioritize by business criticality
			\end{itemize}
		\end{block}
		\begin{alertblock}{Critical Trap}
			Fund remediation by risk reduction, not by age or sunk cost.
		\end{alertblock}
	\end{frame}
	
	\begin{frame}{Q304: Shadow IT Response}
		\begin{block}{Question}
			\small Business units are adopting unsanctioned SaaS tools that process customer data. Blocking them outright would disrupt operations. What is the BEST architecture-aligned response?
		\end{block}
		\begin{enumerate}[(A)]
			\footnotesize
			\item Issue a policy prohibiting all unsanctioned tools and rely on disciplinary action to drive compliance across business units
			\item Discover and assess the tools, bring high-value ones under governance through risk-based onboarding, and offer sanctioned alternatives for unmet needs
			\item Allow continued use but require each business unit to self-certify the security of its tools once a year
			\item Transfer accountability for the tools to the business units that procured them and remove them from the IT risk register
		\end{enumerate}
	\end{frame}
	
	\begin{frame}{Q304: Explanation}
		\begin{block}{Correct Answer: B}
			\textbf{Why (B) is the best answer:}
			\begin{itemize}
				\item Visibility comes before control
				\item Onboard valuable tools through risk-based review
				\item Address the unmet need that drove adoption
			\end{itemize}
		\end{block}
		\begin{alertblock}{Critical Trap}
			Shadow IT is a symptom of unmet needs; prohibition alone drives it underground.
		\end{alertblock}
	\end{frame}
	
	\begin{frame}{Q305: Identity Provider Dependency}
		\begin{block}{Question}
			\small An architecture review finds that authentication for 40 critical applications depends on a single identity provider instance in one data center. What is the MOST appropriate recommendation?
		\end{block}
		\begin{enumerate}[(A)]
			\footnotesize
			\item Migrate all 40 applications to local authentication so that none depend on a shared service for sign-in
			\item Assess the business impact of identity provider failure and design resilience, such as redundancy, failover, and break-glass access, proportionate to application criticality
			\item Accept the risk, because commercial identity platforms are vendor-supported and therefore inherently highly available
			\item Add monitoring to detect outages quickly, since detection alone satisfies the availability requirement of dependent applications
		\end{enumerate}
	\end{frame}
	
	\begin{frame}{Q305: Explanation}
		\begin{block}{Correct Answer: B}
			\textbf{Why (B) is the best answer:}
			\begin{itemize}
				\item A shared dependency creates correlated failure
				\item Resilience should match dependent criticality
				\item Break-glass access protects against lockout
			\end{itemize}
		\end{block}
		\begin{alertblock}{Critical Trap}
			Shared services need resilience sized to what depends on them.
		\end{alertblock}
	\end{frame}
	
	\begin{frame}{Q306: Proprietary Cloud Services}
		\begin{block}{Question}
			\small A proposed cloud platform offers proprietary services that greatly reduce development time but make later migration costly. Which approach BEST addresses the lock-in risk?
		\end{block}
		\begin{enumerate}[(A)]
			\footnotesize
			\item Reject every proprietary service, since lock-in risk always outweighs the benefit of faster delivery
			\item Rely on termination clauses in the contract, since legal exit rights remove the technical dependency
			\item Evaluate exit cost and strategy during the decision, apply abstraction where the business case justifies it, and document residual lock-in as accepted risk
			\item Adopt a multi-cloud design for every workload regardless of criticality to eliminate lock-in entirely
		\end{enumerate}
	\end{frame}
	
	\begin{frame}{Q306: Explanation}
		\begin{block}{Correct Answer: C}
			\textbf{Why (C) is the best answer:}
			\begin{itemize}
				\item Lock-in is a trade-off, not an automatic veto
				\item Plan the exit before committing
				\item Document residual risk for acceptance
			\end{itemize}
		\end{block}
		\begin{alertblock}{Critical Trap}
			Manage lock-in through informed trade-offs, not blanket rules.
		\end{alertblock}
	\end{frame}
	
	\begin{frame}{Q307: Undocumented APIs}
		\begin{block}{Question}
			\small An organization adopting microservices discovers hundreds of undocumented internal APIs. Which risk should the practitioner treat as MOST critical from a governance standpoint?
		\end{block}
		\begin{enumerate}[(A)]
			\footnotesize
			\item Higher network latency from excessive service calls, which will degrade user experience across the platform
			\item Unknown APIs cannot be inventoried, secured, or monitored, which leaves unmanaged attack surface and unclear ownership
			\item Increased licensing cost, because each undocumented API will eventually require its own gateway license
			\item Developers may rename endpoints, which would break backward compatibility for external partners
		\end{enumerate}
	\end{frame}
	
	\begin{frame}{Q307: Explanation}
		\begin{block}{Correct Answer: B}
			\textbf{Why (B) is the best answer:}
			\begin{itemize}
				\item You cannot protect what you cannot see
				\item Ownership of each API must be known
				\item Inventory enables monitoring and security testing
			\end{itemize}
		\end{block}
		\begin{alertblock}{Critical Trap}
			Undocumented interfaces are unmanaged attack surface.
		\end{alertblock}
	\end{frame}
	
	\begin{frame}{Q308: Expired Exception}
		\begin{block}{Question}
			\small An architecture exception allowed a non-standard database for 12 months. The exception has now expired, but the application still runs on it. What should the risk practitioner do FIRST?
		\end{block}
		\begin{enumerate}[(A)]
			\footnotesize
			\item Allow the exception to roll over automatically, since the application remains operational and unchanged
			\item Instruct the application team to migrate immediately, regardless of cost, dependencies, or business impact
			\item Ensure the expired exception is reassessed, current risk is reflected in the risk register, and renewal or remediation is decided by the accountable authority
			\item Close the exception record, since expiry of the approval period shows the matter has been resolved
		\end{enumerate}
	\end{frame}
	
	\begin{frame}{Q308: Explanation}
		\begin{block}{Correct Answer: C}
			\textbf{Why (C) is the best answer:}
			\begin{itemize}
				\item Expiry does not remove the underlying risk
				\item Reassess and update the risk register
				\item Decision authority renews or remediates
			\end{itemize}
		\end{block}
		\begin{alertblock}{Critical Trap}
			An expired exception with a live deviation is an unmanaged risk.
		\end{alertblock}
	\end{frame}
	
	\begin{frame}{Q309: Acquisition Integration}
		\begin{block}{Question}
			\small After an acquisition, leadership wants fast integration of networks and applications. What should be the architecture risk priority?
		\end{block}
		\begin{enumerate}[(A)]
			\footnotesize
			\item Connect the networks immediately to realize synergies, and remediate security issues after integration is complete
			\item Replace all acquired applications with the parent's standard stack before connecting any systems
			\item Rely on the acquired company's compliance attestation, since it was reviewed during due diligence
			\item Assess the acquired environment's security posture and dependencies before connecting it, then integrate in stages with defined trust boundaries
		\end{enumerate}
	\end{frame}
	
	\begin{frame}{Q309: Explanation}
		\begin{block}{Correct Answer: D}
			\textbf{Why (D) is the best answer:}
			\begin{itemize}
				\item Connecting unknown environments imports their risk
				\item Stage integration with trust boundaries
				\item Due diligence is not a substitute for technical assessment
			\end{itemize}
		\end{block}
		\begin{alertblock}{Critical Trap}
			Assess before you connect.
		\end{alertblock}
	\end{frame}
	
	\begin{frame}{Q310: Conflicting Customer Data}
		\begin{block}{Question}
			\small Several departments keep their own copies of customer data with inconsistent definitions, causing reporting errors. Which architecture-level response BEST addresses the root cause?
		\end{block}
		\begin{enumerate}[(A)]
			\footnotesize
			\item Consolidate all departmental databases onto a single server so that duplicate copies disappear physically
			\item Run a nightly reconciliation job that overwrites departmental data with the most recently updated version
			\item Establish authoritative data sources and enterprise definitions with accountable data owners, supported by integration patterns that limit uncontrolled copies
			\item Require each department to document its own definitions in a local data dictionary maintained by its team
		\end{enumerate}
	\end{frame}
	
	\begin{frame}{Q310: Explanation}
		\begin{block}{Correct Answer: C}
			\textbf{Why (C) is the best answer:}
			\begin{itemize}
				\item The root cause is lack of authority and definition
				\item Owners are accountable for each domain
				\item Integration patterns reduce uncontrolled copies
			\end{itemize}
		\end{block}
		\begin{alertblock}{Critical Trap}
			Fix data conflicts with ownership and authoritative sources, not with copying.
		\end{alertblock}
	\end{frame}
	
	\begin{frame}{Q311: Standards and Innovation}
		\begin{block}{Question}
			\small Engineering teams complain that mandatory technology standards slow innovation. What is the BEST way for architecture governance to respond?
		\end{block}
		\begin{enumerate}[(A)]
			\footnotesize
			\item Remove standards for non-production environments, since production controls are sufficient to manage risk
			\item Keep standards for critical controls while providing a defined, risk-based path to evaluate and onboard new technologies
			\item Suspend standards for innovation projects until they reach production readiness, then enforce them
			\item Retain all standards unchanged, since any flexibility inevitably increases risk to the enterprise
		\end{enumerate}
	\end{frame}
	
	\begin{frame}{Q311: Explanation}
		\begin{block}{Correct Answer: B}
			\textbf{Why (B) is the best answer:}
			\begin{itemize}
				\item Governance should enable as well as restrict
				\item Critical controls stay mandatory
				\item A clear onboarding path reduces workarounds
			\end{itemize}
		\end{block}
		\begin{alertblock}{Critical Trap}
			Good governance offers a path for new technology, not only a wall.
		\end{alertblock}
	\end{frame}
	
	\begin{frame}{Q312: Fragile Strategic Capability}
		\begin{block}{Question}
			\small A capability heat map shows a capability rated critical to strategy but supported by an end-of-life application with little investment planned. What should the risk practitioner conclude?
		\end{block}
		\begin{enumerate}[(A)]
			\footnotesize
			\item The risk is low, because end-of-life applications are normally covered by extended support arrangements
			\item The application must be decommissioned immediately, because end-of-life status always outweighs capability value
			\item Strategic dependence on a fragile asset is a significant risk that requires prioritized investment or a mitigation plan
			\item The heat map is flawed, because capabilities should not be rated by strategic importance
		\end{enumerate}
	\end{frame}
	
	\begin{frame}{Q312: Explanation}
		\begin{block}{Correct Answer: C}
			\textbf{Why (C) is the best answer:}
			\begin{itemize}
				\item High criticality plus weak support equals high risk
				\item Mitigation or investment must be planned
				\item Strategy depends on supporting assets
			\end{itemize}
		\end{block}
		\begin{alertblock}{Critical Trap}
			Criticality without support is a risk hotspot.
		\end{alertblock}
	\end{frame}
	
	\begin{frame}{Q313: Reference Architecture Deviation}
		\begin{block}{Question}
			\small A project team deviates from the approved reference architecture to meet a deadline, claiming equivalent controls. What is the BEST way to validate the claim?
		\end{block}
		\begin{enumerate}[(A)]
			\footnotesize
			\item Accept the claim, because the project manager is accountable for delivery and schedule risk
			\item Require full adoption of the reference architecture without considering the control objectives the team cites
			\item Ask the vendor who supplied the components to certify that the design is equivalent
			\item Have architecture or security review the deviation against the control objectives the reference architecture satisfies, and document the outcome
		\end{enumerate}
	\end{frame}
	
	\begin{frame}{Q313: Explanation}
		\begin{block}{Correct Answer: D}
			\textbf{Why (D) is the best answer:}
			\begin{itemize}
				\item Equivalence must be demonstrated, not asserted
				\item Compare against control objectives
				\item Document the review outcome
			\end{itemize}
		\end{block}
		\begin{alertblock}{Critical Trap}
			A claim of equivalent controls needs independent validation.
		\end{alertblock}
	\end{frame}
	
	\begin{frame}{Q314: Cloud-First Strategy}
		\begin{block}{Question}
			\small The CIO announces a ``cloud-first'' strategy. Which action BEST ensures that risk is considered in its implementation?
		\end{block}
		\begin{enumerate}[(A)]
			\footnotesize
			\item Embed risk criteria in workload placement decisions, such as data sensitivity, regulatory constraints, and exit feasibility
			\item Migrate low-risk workloads first and postpone any risk assessment until high-risk workloads are considered
			\item Delegate cloud risk decisions to the provider through the shared responsibility model
			\item Require all business units to adopt cloud by a fixed date to ensure a consistent approach
		\end{enumerate}
	\end{frame}
	
	\begin{frame}{Q314: Explanation}
		\begin{block}{Correct Answer: A}
			\textbf{Why (A) is the best answer:}
			\begin{itemize}
				\item Strategy needs decision criteria
				\item Placement depends on data and regulation
				\item Shared responsibility does not transfer accountability
			\end{itemize}
		\end{block}
		\begin{alertblock}{Critical Trap}
			Cloud-first still needs risk-based placement.
		\end{alertblock}
	\end{frame}
	
	\begin{frame}{Q315: Tight Coupling}
		\begin{block}{Question}
			\small A proposed architecture uses tightly coupled components to reduce cost. What risk should the practitioner highlight?
		\end{block}
		\begin{enumerate}[(A)]
			\footnotesize
			\item Tight coupling increases the likelihood that one failure cascades across components, reducing resilience and making change harder to isolate
			\item Tight coupling prevents scalability, because coupled components can never be deployed on separate servers
			\item Tightly coupled systems cannot be monitored, because telemetry requires loosely coupled interfaces
			\item Cost optimization is not an architecture concern and should not influence design decisions
		\end{enumerate}
	\end{frame}
	
	\begin{frame}{Q315: Explanation}
		\begin{block}{Correct Answer: A}
			\textbf{Why (A) is the best answer:}
			\begin{itemize}
				\item Coupling spreads failures
				\item Changes become harder to isolate
				\item Cost savings must be weighed against resilience
			\end{itemize}
		\end{block}
		\begin{alertblock}{Critical Trap}
			Efficiency gained through coupling is paid for in fragility.
		\end{alertblock}
	\end{frame}
	
	\begin{frame}{Q316: Stale Architecture Repository}
		\begin{block}{Question}
			\small A review finds the enterprise architecture repository is largely out of date. What is the MOST significant implication for risk management?
		\end{block}
		\begin{enumerate}[(A)]
			\footnotesize
			\item Audit findings will increase, since documentation deficiencies are always scored as high severity
			\item Risk assessments and impact analyses based on inaccurate architecture information may misjudge dependencies and exposure
			\item Architecture staff will lose motivation, reducing the quality of future designs
			\item The repository can be disregarded, because operational teams already hold accurate tacit knowledge
		\end{enumerate}
	\end{frame}
	
	\begin{frame}{Q316: Explanation}
		\begin{block}{Correct Answer: B}
			\textbf{Why (B) is the best answer:}
			\begin{itemize}
				\item Risk analysis relies on accurate dependency data
				\item Stale data leads to wrong impact conclusions
				\item Tacit knowledge is not auditable or scalable
			\end{itemize}
		\end{block}
		\begin{alertblock}{Critical Trap}
			Outdated architecture data produces unreliable risk conclusions.
		\end{alertblock}
	\end{frame}
	
	\begin{frame}{Q317: Governance Effectiveness}
		\begin{block}{Question}
			\small How can the risk practitioner BEST determine whether architecture governance is effective?
		\end{block}
		\begin{enumerate}[(A)]
			\footnotesize
			\item Assess whether the process consistently influences decisions, for example exceptions are tracked, deviations are remediated, and outcomes align with principles
			\item Count the number of architecture review meetings held each quarter
			\item Measure the number of projects approved without conditions, since that indicates alignment
			\item Verify that all architecture documents are stored in the repository, whether or not they are used
		\end{enumerate}
	\end{frame}
	
	\begin{frame}{Q317: Explanation}
		\begin{block}{Correct Answer: A}
			\textbf{Why (A) is the best answer:}
			\begin{itemize}
				\item Effectiveness is shown by influence on decisions
				\item Track exceptions and remediation
				\item Activity counts are not outcomes
			\end{itemize}
		\end{block}
		\begin{alertblock}{Critical Trap}
			Measure outcomes of governance, not its activity.
		\end{alertblock}
	\end{frame}
		
	% ============================================================
	\section{2. IT Operations and Change Management}
	% ============================================================
	
	\begin{frame}{Q318: Unapproved Emergency Change}
		\begin{block}{Question}
			\small A critical outage required an emergency change. The change manager was unavailable and an engineer deployed the fix under verbal approval. What is the MOST important follow-up?
		\end{block}
		\begin{enumerate}[(A)]
			\footnotesize
			\item Retroactively review and document the change through the emergency change process, including approval, testing evidence, and post-implementation validation
			\item Take no further action, since the outage was resolved and the speed of response was justified
			\item Revert the change at once to restore the prior state, and then process it through the normal procedure
			\item Discipline the engineer for bypassing approval and treat the matter as closed
		\end{enumerate}
	\end{frame}
	
	\begin{frame}{Q318: Explanation}
		\begin{block}{Correct Answer: A}
			\textbf{Why (A) is the best answer:}
			\begin{itemize}
				\item Emergency changes still need records
				\item Approval and evidence are captured after the event
				\item Validate that the fix did not introduce new risk
			\end{itemize}
		\end{block}
		\begin{alertblock}{Critical Trap}
			Emergency does not mean undocumented.
		\end{alertblock}
	\end{frame}
	
	\begin{frame}{Q319: Rubber-Stamp CAB}
		\begin{block}{Question}
			\small A change advisory board approves 98 percent of changes, while change-related incidents are rising. What should the risk practitioner investigate FIRST?
		\end{block}
		\begin{enumerate}[(A)]
			\footnotesize
			\item Whether the board receives sufficient risk and impact information and challenges changes rigorously, or approves as a formality
			\item Whether the board should meet more often to reduce its backlog of pending requests
			\item Whether membership should be limited to IT staff so that decisions are faster
			\item Whether incidents are being misclassified, since a high approval rate indicates good change quality
		\end{enumerate}
	\end{frame}
	
	\begin{frame}{Q319: Explanation}
		\begin{block}{Correct Answer: A}
			\textbf{Why (A) is the best answer:}
			\begin{itemize}
				\item Approval rate alone says little
				\item Look at the quality of information reviewed
				\item Challenge and testing evidence matter
			\end{itemize}
		\end{block}
		\begin{alertblock}{Critical Trap}
			A board that approves everything is not a control.
		\end{alertblock}
	\end{frame}
	
	\begin{frame}{Q320: Developer Deploys to Production}
		\begin{block}{Question}
			\small A developer can deploy their own code to production in an organization that requires change approval. If segregation cannot be fully achieved, which is the BEST compensating control?
		\end{block}
		\begin{enumerate}[(A)]
			\footnotesize
			\item A manager's informal confirmation that the developer is trustworthy and experienced
			\item Independent review of deployment logs against approved change records, performed by someone other than the deployer
			\item Quarterly security awareness training that covers secure deployment practices
			\item A policy statement reminding developers that unapproved code must not be deployed
		\end{enumerate}
	\end{frame}
	
	\begin{frame}{Q320: Explanation}
		\begin{block}{Correct Answer: B}
			\textbf{Why (B) is the best answer:}
			\begin{itemize}
				\item Compensating controls must detect misuse
				\item Independence is essential
				\item Reconcile deployments to approvals
			\end{itemize}
		\end{block}
		\begin{alertblock}{Critical Trap}
			A compensating control needs independent detection, not trust.
		\end{alertblock}
	\end{frame}
	
	\begin{frame}{Q321: Inaccurate CMDB}
		\begin{block}{Question}
			\small An audit finds the CMDB is accurate for only 60 percent of critical assets. What is the MOST significant risk?
		\end{block}
		\begin{enumerate}[(A)]
			\footnotesize
			\item Asset depreciation will be misreported in the financial statements
			\item Software licensing costs will rise, because vendors can audit the CMDB for discrepancies
			\item Incident, change, and vulnerability decisions will rely on incomplete dependency and ownership data, causing misjudged impact
			\item IT staff will be unable to access configuration management tools during incidents
		\end{enumerate}
	\end{frame}
	
	\begin{frame}{Q321: Explanation}
		\begin{block}{Correct Answer: C}
			\textbf{Why (C) is the best answer:}
			\begin{itemize}
				\item CMDB supports many operational processes
				\item Poor data misleads impact assessment
				\item Ownership gaps delay response
			\end{itemize}
		\end{block}
		\begin{alertblock}{Critical Trap}
			A weak CMDB weakens every process that depends on it.
		\end{alertblock}
	\end{frame}
	
	\begin{frame}{Q322: Unpatchable Legacy System}
		\begin{block}{Question}
			\small A business-critical system cannot be patched because the vendor no longer supports it, and a high-severity vulnerability is being exploited in the wild. What is the BEST response?
		\end{block}
		\begin{enumerate}[(A)]
			\footnotesize
			\item Accept the vulnerability, because unsupported systems cannot be remediated in any way
			\item Implement compensating controls such as isolation, access restriction, and enhanced monitoring, document risk acceptance, and plan replacement
			\item Shut down the system immediately regardless of its business impact
			\item Apply a patch from a public forum that claims to fix the flaw, since it is the only option available
		\end{enumerate}
	\end{frame}
	
	\begin{frame}{Q322: Explanation}
		\begin{block}{Correct Answer: B}
			\textbf{Why (B) is the best answer:}
			\begin{itemize}
				\item Isolate and monitor the exposed system
				\item Document acceptance by the accountable owner
				\item Plan replacement of unsupported technology
			\end{itemize}
		\end{block}
		\begin{alertblock}{Critical Trap}
			If you cannot patch, you must contain.
		\end{alertblock}
	\end{frame}
	
	\begin{frame}{Q323: Recurring Incidents}
		\begin{block}{Question}
			\small Incidents of the same type recur monthly and are each resolved quickly by the service desk. What does this MOST likely indicate?
		\end{block}
		\begin{enumerate}[(A)]
			\footnotesize
			\item Problem management is ineffective, because root causes are not being identified and eliminated
			\item Incident management is working well, since each incident is resolved within the SLA
			\item The service desk needs additional staff to reduce recurrence
			\item The incident classification scheme should be simplified to reduce repeated tickets
		\end{enumerate}
	\end{frame}
	
	\begin{frame}{Q323: Explanation}
		\begin{block}{Correct Answer: A}
			\textbf{Why (A) is the best answer:}
			\begin{itemize}
				\item Fast fixes can hide systemic causes
				\item Problem management targets root cause
				\item Recurrence signals a process gap
			\end{itemize}
		\end{block}
		\begin{alertblock}{Critical Trap}
			Quick resolution of repeat incidents hides a problem-management failure.
		\end{alertblock}
	\end{frame}
	
	\begin{frame}{Q324: SLA Met but Users Unhappy}
		\begin{block}{Question}
			\small A provider consistently meets its SLA availability targets, yet business users complain about poor performance. What is the BEST explanation?
		\end{block}
		\begin{enumerate}[(A)]
			\footnotesize
			\item SLA metrics may not reflect the business-relevant user experience, such as measuring server uptime rather than transaction performance
			\item Users are exaggerating, since SLA metrics are objective and verified
			\item The SLA is too strict and should be relaxed to match what users realistically expect
			\item The provider should be penalized regardless, because complaints indicate a breach
		\end{enumerate}
	\end{frame}
	
	\begin{frame}{Q324: Explanation}
		\begin{block}{Correct Answer: A}
			\textbf{Why (A) is the best answer:}
			\begin{itemize}
				\item Metrics must reflect business outcomes
				\item Uptime is not the same as usable service
				\item Revisit SLA measures with the business
			\end{itemize}
		\end{block}
		\begin{alertblock}{Critical Trap}
			Meeting the SLA is not the same as meeting the need.
		\end{alertblock}
	\end{frame}
	
	\begin{frame}{Q325: Peak Season Failure}
		\begin{block}{Question}
			\small A retailer's e-commerce platform failed in peak season despite monthly capacity reports showing 60 percent average utilization. What is the MOST likely gap?
		\end{block}
		\begin{enumerate}[(A)]
			\footnotesize
			\item Capacity planning based on average utilization did not model peak demand or growth scenarios
			\item Capacity reports should have been reviewed weekly rather than monthly
			\item Servers should have been sized to run at 100 percent utilization targets
			\item Capacity planning is the cloud provider's responsibility under any deployment model
		\end{enumerate}
	\end{frame}
	
	\begin{frame}{Q325: Explanation}
		\begin{block}{Correct Answer: A}
			\textbf{Why (A) is the best answer:}
			\begin{itemize}
				\item Averages hide peaks
				\item Model growth and seasonal scenarios
				\item Plan for demand, not for history
			\end{itemize}
		\end{block}
		\begin{alertblock}{Critical Trap}
			Plan capacity for peaks, not averages.
		\end{alertblock}
	\end{frame}
	
	\begin{frame}{Q326: Untested Rollback}
		\begin{block}{Question}
			\small A release process has no tested rollback procedure. A deployment fails and takes hours to reverse. Which control improvement is MOST appropriate?
		\end{block}
		\begin{enumerate}[(A)]
			\footnotesize
			\item Prohibit deployments during business hours so failures affect fewer users
			\item Require documented and tested rollback plans as an approval criterion, validated in pre-production
			\item Increase the number of approvers required on every release
			\item Require all releases to be deployed by the most senior engineers
		\end{enumerate}
	\end{frame}
	
	\begin{frame}{Q326: Explanation}
		\begin{block}{Correct Answer: B}
			\textbf{Why (B) is the best answer:}
			\begin{itemize}
				\item Rollback is a recovery control for change
				\item It must be tested before approval
				\item Approval criteria enforce discipline
			\end{itemize}
		\end{block}
		\begin{alertblock}{Critical Trap}
			An untested rollback plan is only a hope.
		\end{alertblock}
	\end{frame}
	
	\begin{frame}{Q327: Alert Fatigue}
		\begin{block}{Question}
			\small Security operations receive thousands of alerts daily and analysts routinely ignore them. What is the BEST way to improve detection effectiveness?
		\end{block}
		\begin{enumerate}[(A)]
			\footnotesize
			\item Disable all low-severity alerts across every system to cut volume
			\item Hire additional analysts so that every alert can be reviewed manually
			\item Outsource alert review to a provider without changing the existing alert rules
			\item Tune and prioritize alerts using risk-based use cases, thresholds, and context so analysts focus on high-fidelity events
		\end{enumerate}
	\end{frame}
	
	\begin{frame}{Q327: Explanation}
		\begin{block}{Correct Answer: D}
			\textbf{Why (D) is the best answer:}
			\begin{itemize}
				\item Noise hides real incidents
				\item Tune with risk-based use cases
				\item Add context to improve fidelity
			\end{itemize}
		\end{block}
		\begin{alertblock}{Critical Trap}
			More alerts are not better detection.
		\end{alertblock}
	\end{frame}
	
	\begin{frame}{Q328: Shared Admin Accounts}
		\begin{block}{Question}
			\small A managed service provider administers production systems using shared privileged accounts. What is the PRIMARY risk?
		\end{block}
		\begin{enumerate}[(A)]
			\footnotesize
			\item Password complexity cannot be enforced on shared accounts, which makes brute-force attacks the primary exposure from this arrangement
			\item Lack of individual accountability and inability to attribute actions, increasing the risk of undetected misuse
			\item Provider staff will be unable to perform emergency changes, because shared accounts are blocked by most change management tools
			\item Contractual liability will shift to the organization, because shared credential use is treated as customer-approved access under standard service terms
		\end{enumerate}
	\end{frame}
	
	\begin{frame}{Q328: Explanation}
		\begin{block}{Correct Answer: B}
			\textbf{Why (B) is the best answer:}
			\begin{itemize}
				\item Shared accounts prevent attribution
				\item Misuse becomes hard to detect
				\item Use named accounts with PAM
			\end{itemize}
		\end{block}
		\begin{alertblock}{Critical Trap}
			Shared credentials destroy accountability.
		\end{alertblock}
	\end{frame}
	
	\begin{frame}{Q329: Reliable Asset Inventory}
		\begin{block}{Question}
			\small Which is the MOST reliable approach to maintaining an accurate hardware and software inventory?
		\end{block}
		\begin{enumerate}[(A)]
			\footnotesize
			\item Conduct an annual manual count performed by IT staff
			\item Ask department heads to email lists of the assets they use
			\item Rely on purchase orders alone as the source of truth
			\item Combine automated discovery with reconciliation to procurement and ownership records, and investigate discrepancies
		\end{enumerate}
	\end{frame}
	
	\begin{frame}{Q329: Explanation}
		\begin{block}{Correct Answer: D}
			\textbf{Why (D) is the best answer:}
			\begin{itemize}
				\item Automated discovery finds the unknown
				\item Reconcile with procurement and owners
				\item Discrepancies reveal shadow assets
			\end{itemize}
		\end{block}
		\begin{alertblock}{Critical Trap}
			Inventory accuracy needs discovery plus reconciliation.
		\end{alertblock}
	\end{frame}
	
	\begin{frame}{Q330: Pre-Approved Standard Changes}
		\begin{block}{Question}
			\small A DevOps team wants to pre-approve frequent low-risk deployments to avoid CAB delays. What is the BEST risk-informed approach?
		\end{block}
		\begin{enumerate}[(A)]
			\footnotesize
			\item Exempt DevOps teams from change management entirely, since they own their deployments and are accountable for the results
			\item Define standard changes with clear criteria, automated testing and controls, and periodic review of outcomes to confirm they remain low risk
			\item Pre-approve all changes below a defined cost threshold, since low-cost changes carry proportionately low risk
			\item Allow teams to classify any change as standard without documented criteria, provided the team lead agrees afterward
		\end{enumerate}
	\end{frame}
	
	\begin{frame}{Q330: Explanation}
		\begin{block}{Correct Answer: B}
			\textbf{Why (B) is the best answer:}
			\begin{itemize}
				\item Pre-approval needs defined criteria
				\item Automation provides control evidence
				\item Review outcomes to validate low risk
			\end{itemize}
		\end{block}
		\begin{alertblock}{Critical Trap}
			Speed is gained through controlled standardization, not exemption.
		\end{alertblock}
	\end{frame}
	
	\begin{frame}{Q331: Skipped Post-Implementation Reviews}
		\begin{block}{Question}
			\small Post-implementation reviews are skipped for most changes. What is the BEST justification for reinstating them?
		\end{block}
		\begin{enumerate}[(A)]
			\footnotesize
			\item They verify that intended outcomes and controls were achieved, capture lessons learned, and reveal unexpected or unauthorized side effects
			\item They satisfy a regulatory requirement that applies to all changes without exception
			\item They allow IT leadership to assign accountability for failed changes and apply performance penalties to the responsible teams
			\item They replace the need for testing before deployment by validating each change in the live environment instead
		\end{enumerate}
	\end{frame}
	
	\begin{frame}{Q331: Explanation}
		\begin{block}{Correct Answer: A}
			\textbf{Why (A) is the best answer:}
			\begin{itemize}
				\item Confirm objectives and controls were met
				\item Capture lessons for future changes
				\item Detect unintended effects
			\end{itemize}
		\end{block}
		\begin{alertblock}{Critical Trap}
			Reviews close the change loop.
		\end{alertblock}
	\end{frame}
	
	\begin{frame}{Q332: Monitoring Coverage}
		\begin{block}{Question}
			\small Which situation represents the greatest operational risk related to monitoring?
		\end{block}
		\begin{enumerate}[(A)]
			\footnotesize
			\item Monitoring dashboards differ in layout and color across teams, which prevents consistent interpretation of system status during incidents
			\item Monitoring data is retained longer than operational needs require, increasing storage cost beyond what the business case approved
			\item Non-production systems are not monitored in real time, which leaves development environments without immediate failure alerts
			\item Critical systems are monitored for availability but not for security-relevant events or performance degradation
		\end{enumerate}
	\end{frame}
	
	\begin{frame}{Q332: Explanation}
		\begin{block}{Correct Answer: D}
			\textbf{Why (D) is the best answer:}
			\begin{itemize}
				\item Coverage must match risk
				\item Availability alone misses degradation
				\item Security events need visibility
			\end{itemize}
		\end{block}
		\begin{alertblock}{Critical Trap}
			Monitor what matters, not just what is easy.
		\end{alertblock}
	\end{frame}
	
	\begin{frame}{Q333: Silent Batch Failure}
		\begin{block}{Question}
			\small A batch job failed silently over a weekend, causing incorrect financial reports. Which control improvement is MOST appropriate?
		\end{block}
		\begin{enumerate}[(A)]
			\footnotesize
			\item Run batch jobs only during business hours so failures are noticed sooner
			\item Require the CFO to manually review all report outputs
			\item Replace the scheduling tool with a different vendor's product
			\item Implement job monitoring with automated failure alerting, escalation, and reconciliation checks on outputs
		\end{enumerate}
	\end{frame}
	
	\begin{frame}{Q333: Explanation}
		\begin{block}{Correct Answer: D}
			\textbf{Why (D) is the best answer:}
			\begin{itemize}
				\item Failures must generate alerts
				\item Escalation ensures response
				\item Reconcile outputs to detect errors
			\end{itemize}
		\end{block}
		\begin{alertblock}{Critical Trap}
			A failure nobody sees is a control gap.
		\end{alertblock}
	\end{frame}
	
	\begin{frame}{Q334: Missed Escalation}
		\begin{block}{Question}
			\small A major incident exceeded its expected resolution time, but the incident manager never escalated to senior management. What is the MOST likely root cause to examine?
		\end{block}
		\begin{enumerate}[(A)]
			\footnotesize
			\item Insufficient technical skill in the resolver team, which delayed diagnosis and made escalation seem unnecessary
			\item Excessive reliance on automated triage tools, which suppressed the alerts that would otherwise have prompted escalation
			\item Lack of interest in operational issues among senior management, which discouraged the incident manager from raising concerns
			\item Absence of clear escalation criteria and thresholds based on impact and duration in the incident management process
		\end{enumerate}
	\end{frame}
	
	\begin{frame}{Q334: Explanation}
		\begin{block}{Correct Answer: D}
			\textbf{Why (D) is the best answer:}
			\begin{itemize}
				\item Escalation needs defined triggers
				\item Impact and duration drive thresholds
				\item Process gaps beat blaming individuals
			\end{itemize}
		\end{block}
		\begin{alertblock}{Critical Trap}
			Without triggers, escalation depends on judgment under pressure.
		\end{alertblock}
	\end{frame}
	
	\begin{frame}{Q335: Cloud Console Changes}
		\begin{block}{Question}
			\small Operations teams lack visibility into configuration changes made through cloud management consoles. What is the BEST control?
		\end{block}
		\begin{enumerate}[(A)]
			\footnotesize
			\item Permanently disable console access for all staff and require every change to go through a ticket handled by a single administrator
			\item Rely on provider-level logs for audit purposes, without a defined process to review or alert on them
			\item Enforce infrastructure as code with change logging, and detect and alert on drift from approved configuration
			\item Require staff to notify the team by email after each console change so that changes can be recorded manually
		\end{enumerate}
	\end{frame}
	
	\begin{frame}{Q335: Explanation}
		\begin{block}{Correct Answer: C}
			\textbf{Why (C) is the best answer:}
			\begin{itemize}
				\item Code-driven changes are reviewable
				\item Drift detection finds unapproved changes
				\item Logs must be monitored, not just kept
			\end{itemize}
		\end{block}
		\begin{alertblock}{Critical Trap}
			Controlled cloud change comes from code and drift detection.
		\end{alertblock}
	\end{frame}
		
	% ============================================================
	\section{3. Enterprise Resiliency}
	% ============================================================
	
	\begin{frame}{Q336: RTO Versus Capability}
		\begin{block}{Question}
			\small A BIA sets an RTO of 4 hours for a process, but the supporting application has a documented recovery time of 12 hours. What is the BEST action?
		\end{block}
		\begin{enumerate}[(A)]
			\footnotesize
			\item Address the gap by improving recovery capability or having the business formally reassess and accept the risk based on impact
			\item Update the BIA to reflect the 12-hour recovery time as the new RTO
			\item Take no action, since a DR plan exists and is documented
			\item Move the process to a different application without further analysis
		\end{enumerate}
	\end{frame}
	
	\begin{frame}{Q336: Explanation}
		\begin{block}{Correct Answer: A}
			\textbf{Why (A) is the best answer:}
			\begin{itemize}
				\item A gap between need and capability is a risk
				\item Fix capability or formally accept risk
				\item Do not rewrite requirements to fit capability
			\end{itemize}
		\end{block}
		\begin{alertblock}{Critical Trap}
			Do not change the requirement to hide the gap.
		\end{alertblock}
	\end{frame}
	
	\begin{frame}{Q337: RPO Mismatch}
		\begin{block}{Question}
			\small A system replicates data to the DR site every 6 hours, while the business states it can tolerate at most 1 hour of data loss. Which statement is correct?
		\end{block}
		\begin{enumerate}[(A)]
			\footnotesize
			\item The RTO is violated because replication occurs too infrequently
			\item The design is acceptable, because RPO applies only to tape backups
			\item The replication interval does not meet the RPO, so the solution must be redesigned or the shortfall formally accepted
			\item The MTPD must be extended to 6 hours to match the replication interval
		\end{enumerate}
	\end{frame}
	
	\begin{frame}{Q337: Explanation}
		\begin{block}{Correct Answer: C}
			\textbf{Why (C) is the best answer:}
			\begin{itemize}
				\item RPO defines tolerable data loss
				\item Replication interval must be at or below the RPO
				\item Shortfall requires a decision
			\end{itemize}
		\end{block}
		\begin{alertblock}{Critical Trap}
			RPO is about data loss, RTO is about time.
		\end{alertblock}
	\end{frame}
	
	\begin{frame}{Q338: RTO and MTPD}
		\begin{block}{Question}
			\small Which relationship between RTO and MTPD is MOST appropriate?
		\end{block}
		\begin{enumerate}[(A)]
			\footnotesize
			\item RTO is shorter than MTPD, leaving margin before unacceptable impact occurs
			\item RTO equals MTPD, because any margin implies wasted recovery investment
			\item RTO exceeds MTPD, because RTO measures only technical recovery effort
			\item They are unrelated, because they are defined for different audiences
		\end{enumerate}
	\end{frame}
	
	\begin{frame}{Q338: Explanation}
		\begin{block}{Correct Answer: A}
			\textbf{Why (A) is the best answer:}
			\begin{itemize}
				\item MTPD is the limit of tolerable outage
				\item RTO is the target within that limit
				\item Margin allows for recovery variability
			\end{itemize}
		\end{block}
		\begin{alertblock}{Critical Trap}
			RTO must fall inside MTPD.
		\end{alertblock}
	\end{frame}
	
	\begin{frame}{Q339: Unrealistic DR Test}
		\begin{block}{Question}
			\small A DR test was announced two weeks in advance, used the best staff, and restored only non-critical systems. What is the MOST significant limitation?
		\end{block}
		\begin{enumerate}[(A)]
			\footnotesize
			\item The test was too short to evaluate the recovery procedures properly
			\item The test does not give reliable assurance that critical systems can be recovered under realistic conditions
			\item Using the best staff violates segregation of duties requirements
			\item Advance announcement of a DR test is a regulatory violation
		\end{enumerate}
	\end{frame}
	
	\begin{frame}{Q339: Explanation}
		\begin{block}{Correct Answer: B}
			\textbf{Why (B) is the best answer:}
			\begin{itemize}
				\item Tests must reflect real conditions
				\item Critical systems must be included
				\item Surprise and varied staffing build confidence
			\end{itemize}
		\end{block}
		\begin{alertblock}{Critical Trap}
			A test that avoids the hard parts proves little.
		\end{alertblock}
	\end{frame}
	
	\begin{frame}{Q340: Ransomware and Backups}
		\begin{block}{Question}
			\small Following a ransomware event, backups were also encrypted because they were reachable from the production domain. Which control BEST addresses this?
		\end{block}
		\begin{enumerate}[(A)]
			\footnotesize
			\item Immutable or offline backups with separate credentials and tested restoration
			\item Increase backup frequency to every hour, writing to the same repository so that recovery points are as recent as possible
			\item Apply stronger encryption to the backup files so that attackers cannot read the data even if they can reach it
			\item Extend backup retention to twelve months so that older copies exist if recent copies are affected
		\end{enumerate}
	\end{frame}
	
	\begin{frame}{Q340: Explanation}
		\begin{block}{Correct Answer: A}
			\textbf{Why (A) is the best answer:}
			\begin{itemize}
				\item Backups must be isolated from production
				\item Immutability prevents tampering
				\item Restoration must be tested
			\end{itemize}
		\end{block}
		\begin{alertblock}{Critical Trap}
			Backups reachable by attackers are not a recovery plan.
		\end{alertblock}
	\end{frame}
	
	\begin{frame}{Q341: Single Region SLA}
		\begin{block}{Question}
			\small A critical SaaS application is hosted in a single cloud region, and the provider's SLA guarantees 99.9 percent availability. What is the BEST risk consideration?
		\end{block}
		\begin{enumerate}[(A)]
			\footnotesize
			\item SLA credits do not restore operations, so the organization must assess whether regional failure fits within its impact tolerance and plan accordingly
			\item The SLA transfers all availability risk to the provider
			\item A 99.9 percent SLA is equivalent to no unplanned downtime
			\item A single region is always more resilient than multi-region
		\end{enumerate}
	\end{frame}
	
	\begin{frame}{Q341: Explanation}
		\begin{block}{Correct Answer: A}
			\textbf{Why (A) is the best answer:}
			\begin{itemize}
				\item SLA compensates but does not recover
				\item Compare outage risk with tolerance
				\item Plan for regional failure where needed
			\end{itemize}
		\end{block}
		\begin{alertblock}{Critical Trap}
			SLAs compensate for outages; they do not prevent them.
		\end{alertblock}
	\end{frame}
	
	\begin{frame}{Q342: Supplier Without BCP}
		\begin{block}{Question}
			\small A business continuity plan assumes the availability of a critical supplier that has no continuity plan of its own. What is the BEST response?
		\end{block}
		\begin{enumerate}[(A)]
			\footnotesize
			\item Remove the dependency from the plan and document it as out of scope, since the supplier's resilience is outside the organization's control
			\item Accept the supplier's verbal assurance of continuity, since the supplier has an established reputation and a long relationship with the organization
			\item Transfer the risk to the supplier through contract wording alone, and consider the matter closed once the contract is signed
			\item Assess supplier resilience, require or evaluate continuity capability, and develop alternatives or contingencies for the dependency
		\end{enumerate}
	\end{frame}
	
	\begin{frame}{Q342: Explanation}
		\begin{block}{Correct Answer: D}
			\textbf{Why (D) is the best answer:}
			\begin{itemize}
				\item Dependencies must be assessed
				\item Develop contingencies
				\item Contracts do not guarantee recovery
			\end{itemize}
		\end{block}
		\begin{alertblock}{Critical Trap}
			A dependency you cannot verify is a risk you own.
		\end{alertblock}
	\end{frame}
	
	\begin{frame}{Q343: Communication Failure in Exercise}
		\begin{block}{Question}
			\small In a tabletop exercise, technical recovery went well, but communications to customers and regulators were delayed and inconsistent. What is the MOST appropriate improvement?
		\end{block}
		\begin{enumerate}[(A)]
			\footnotesize
			\item Add more technical staff to the response team so that recovery work is completed sooner and communication can follow
			\item Remove external communications from the plan to avoid inconsistent messaging, and handle them ad hoc through executives
			\item Reserve all communications for legal counsel without defined timing, so that every statement is reviewed before release
			\item Define communication roles, templates, and approval paths in the plan and rehearse them
		\end{enumerate}
	\end{frame}
	
	\begin{frame}{Q343: Explanation}
		\begin{block}{Correct Answer: D}
			\textbf{Why (D) is the best answer:}
			\begin{itemize}
				\item Communication is part of recovery
				\item Pre-approved templates save time
				\item Rehearsal exposes delays
			\end{itemize}
		\end{block}
		\begin{alertblock}{Critical Trap}
			Recovery includes communication.
		\end{alertblock}
	\end{frame}
	
	\begin{frame}{Q344: Outdated BCP}
		\begin{block}{Question}
			\small A BCP was last updated three years ago, although significant organizational and technology changes have occurred. What is the PRIMARY risk?
		\end{block}
		\begin{enumerate}[(A)]
			\footnotesize
			\item The plan may no longer reflect critical processes, dependencies, and contacts, making it ineffective in a real event
			\item Auditors will impose a financial penalty, because regulations require continuity plans to be rewritten every year
			\item The plan will become too detailed for executives to use, reducing the likelihood that decision makers will read it
			\item Insurance coverage will lapse automatically, because insurers require evidence of annual plan revisions
		\end{enumerate}
	\end{frame}
	
	\begin{frame}{Q344: Explanation}
		\begin{block}{Correct Answer: A}
			\textbf{Why (A) is the best answer:}
			\begin{itemize}
				\item Plans must follow organizational change
				\item Contacts and dependencies change
				\item Review on a schedule and after major changes
			\end{itemize}
		\end{block}
		\begin{alertblock}{Critical Trap}
			A plan that does not match reality will fail when needed.
		\end{alertblock}
	\end{frame}
	
	\begin{frame}{Q345: Recovery Sequencing}
		\begin{block}{Question}
			\small During DR planning, several applications claim to be Tier 1. What should determine recovery sequencing?
		\end{block}
		\begin{enumerate}[(A)]
			\footnotesize
			\item The seniority and influence of each application owner, since senior owners are best placed to judge their own priority
			\item The technical complexity of each application, so that the simplest systems are restored first to show quick progress
			\item BIA results, including dependencies between applications and shared infrastructure
			\item The order in which the applications were deployed, because older systems usually underpin newer ones
		\end{enumerate}
	\end{frame}
	
	\begin{frame}{Q345: Explanation}
		\begin{block}{Correct Answer: C}
			\textbf{Why (C) is the best answer:}
			\begin{itemize}
				\item Business impact sets priority
				\item Dependencies determine order
				\item Foundation services recover first
			\end{itemize}
		\end{block}
		\begin{alertblock}{Critical Trap}
			Sequence by business impact and dependency, not by politics.
		\end{alertblock}
	\end{frame}
	
	\begin{frame}{Q346: Untested Restores}
		\begin{block}{Question}
			\small Backups complete successfully every night, but restoration has never been tested. What is the BEST assessment?
		\end{block}
		\begin{enumerate}[(A)]
			\footnotesize
			\item The control is effective, because job completion is verified daily and failures are corrected by the backup team
			\item Restoration testing is optional where backups are encrypted, since encryption already assures integrity and availability
			\item Backup completion does not demonstrate recoverability, so periodic restoration tests are required
			\item Testing can be replaced by increasing backup frequency, since more recovery points reduce the need to verify restores
		\end{enumerate}
	\end{frame}
	
	\begin{frame}{Q346: Explanation}
		\begin{block}{Correct Answer: C}
			\textbf{Why (C) is the best answer:}
			\begin{itemize}
				\item Success status is not proof of restore
				\item Test restores regularly
				\item Frequency does not equal recoverability
			\end{itemize}
		\end{block}
		\begin{alertblock}{Critical Trap}
			A backup is only as good as its last successful restore.
		\end{alertblock}
	\end{frame}
	
	\begin{frame}{Q347: Active-Active Assumption}
		\begin{block}{Question}
			\small Management proposes active-active data centers to guarantee zero downtime. Which risk should the practitioner highlight?
		\end{block}
		\begin{enumerate}[(A)]
			\footnotesize
			\item Active-active designs cannot meet recovery time objectives, because switching traffic between sites always takes longer than restoring from backup
			\item Shared logical failures, such as bad deployments, data corruption, or misconfiguration, can propagate to both sites simultaneously
			\item Cost makes active-active unsuitable for every critical application, so asynchronous replication is always the better choice
			\item Both sites will always be affected by the same physical disaster, because active-active requires them to be in the same metropolitan area
		\end{enumerate}
	\end{frame}
	
	\begin{frame}{Q347: Explanation}
		\begin{block}{Correct Answer: B}
			\textbf{Why (B) is the best answer:}
			\begin{itemize}
				\item Replication copies mistakes too
				\item Logical failures affect both sites
				\item Backups and change control are still needed
			\end{itemize}
		\end{block}
		\begin{alertblock}{Critical Trap}
			Redundancy protects against hardware failure, not against bad changes.
		\end{alertblock}
	\end{frame}
	
	\begin{frame}{Q348: Concentration Risk}
		\begin{block}{Question}
			\small Many critical functions rely on a single cloud provider and a single DNS service. Which term BEST describes this risk?
		\end{block}
		\begin{enumerate}[(A)]
			\footnotesize
			\item Concentration risk, where one provider failure could disrupt multiple critical services at once
			\item Residual risk, because the dependency remains after existing controls and redundancy have been applied
			\item Inherent risk, because the dependency exists before any controls are considered and is determined by the business model
			\item Risk appetite, because the board has chosen to rely on a small number of providers to simplify management
		\end{enumerate}
	\end{frame}
	
	\begin{frame}{Q348: Explanation}
		\begin{block}{Correct Answer: A}
			\textbf{Why (A) is the best answer:}
			\begin{itemize}
				\item Many services share one dependency
				\item A single failure creates broad impact
				\item Diversify or build contingency
			\end{itemize}
		\end{block}
		\begin{alertblock}{Critical Trap}
			One dependency, many failures.
		\end{alertblock}
	\end{frame}
	
	\begin{frame}{Q349: Impact Tolerance}
		\begin{block}{Question}
			\small Regulators require firms to define impact tolerances for important business services. How does impact tolerance differ from RTO?
		\end{block}
		\begin{enumerate}[(A)]
			\footnotesize
			\item Impact tolerance is a technical metric measured by IT, expressing how quickly infrastructure must be restored after failure
			\item Impact tolerance replaces the BIA, because it gives a single maximum disruption figure for the entire organization
			\item Impact tolerance is determined by service vendors in the SLA, which sets the limit of disruption customers must accept
			\item Impact tolerance defines the maximum disruption of a business service the organization can withstand, independent of the technical recovery design
		\end{enumerate}
	\end{frame}
	
	\begin{frame}{Q349: Explanation}
		\begin{block}{Correct Answer: D}
			\textbf{Why (D) is the best answer:}
			\begin{itemize}
				\item Tolerance is business-service focused
				\item It sets the limit of acceptable harm
				\item Technical design must meet it
			\end{itemize}
		\end{block}
		\begin{alertblock}{Critical Trap}
			Impact tolerance is a business limit, not a technical target.
		\end{alertblock}
	\end{frame}
	
	\begin{frame}{Q350: People and Site Dependencies}
		\begin{block}{Question}
			\small A BCP focuses on IT recovery but assumes staff can reach alternate sites. A regional event prevents travel. What gap does this show?
		\end{block}
		\begin{enumerate}[(A)]
			\footnotesize
			\item Inadequate backup frequency, which prevented timely restoration of data at the alternate site
			\item Insufficient firewall capacity at the alternate site, which slowed access after failover
			\item Insufficient consideration of people and workspace dependencies and remote work capability
			\item Excessive involvement of vendors in the response, which reduced the organization's control over recovery decisions
		\end{enumerate}
	\end{frame}
	
	\begin{frame}{Q350: Explanation}
		\begin{block}{Correct Answer: C}
			\textbf{Why (C) is the best answer:}
			\begin{itemize}
				\item Recovery depends on people
				\item Alternate access must be tested
				\item Plan beyond technology
			\end{itemize}
		\end{block}
		\begin{alertblock}{Critical Trap}
			Continuity includes people and places, not just systems.
		\end{alertblock}
	\end{frame}
	
	\begin{frame}{Q351: Restoring After Ransomware}
		\begin{block}{Question}
			\small After a ransomware attack, the team wants to restore the latest backup immediately. What should be done FIRST?
		\end{block}
		\begin{enumerate}[(A)]
			\footnotesize
			\item Restore the latest backup immediately to minimize downtime, and run the investigation in parallel once systems are back
			\item Pay the ransom to obtain the decryption key quickly, since this is usually faster than restoring from backups
			\item Rebuild all systems from scratch before confirming scope, so that nothing from the compromised environment remains
			\item Contain the threat and identify a clean restore point to avoid reintroducing the malware
		\end{enumerate}
	\end{frame}
	
	\begin{frame}{Q351: Explanation}
		\begin{block}{Correct Answer: D}
			\textbf{Why (D) is the best answer:}
			\begin{itemize}
				\item Contain before restoring
				\item Choose a verified clean restore point
				\item Restoring infected data re-infects
			\end{itemize}
		\end{block}
		\begin{alertblock}{Critical Trap}
			Speed without containment repeats the incident.
		\end{alertblock}
	\end{frame}
	
	\begin{frame}{Q352: DR Test Frequency}
		\begin{block}{Question}
			\small Which factor should MOST influence the frequency and depth of DR testing for a system?
		\end{block}
		\begin{enumerate}[(A)]
			\footnotesize
			\item The cost of the DR site, since expensive recovery environments should be tested less often to control operating expense
			\item The vendor's recommended schedule, since the vendor best understands how its own recovery features should be verified
			\item Criticality of supported processes, rate of change in the environment, and previous test results
			\item The number of users with access to the system, since a larger user base creates more recovery complexity
		\end{enumerate}
	\end{frame}
	
	\begin{frame}{Q352: Explanation}
		\begin{block}{Correct Answer: C}
			\textbf{Why (C) is the best answer:}
			\begin{itemize}
				\item Test more where impact is greater
				\item Change increases test need
				\item Past failures drive deeper testing
			\end{itemize}
		\end{block}
		\begin{alertblock}{Critical Trap}
			Test in proportion to risk.
		\end{alertblock}
	\end{frame}
		
	% ============================================================
	\section{4. Data Lifecycle Management}
	% ============================================================
	
	\begin{frame}{Q353: Overcomplicated Classification}
		\begin{block}{Question}
			\small Data classification is applied inconsistently because users find a seven-level scheme too complex. What is the BEST improvement?
		\end{block}
		\begin{enumerate}[(A)]
			\footnotesize
			\item Mandate disciplinary action for any misclassification, so that users treat the scheme seriously
			\item Increase the number of levels so that each data type can be classified with greater precision
			\item Simplify the scheme to a few clear levels with handling rules and examples, supported by training and automated labeling
			\item Classify all data at the highest level so that users do not need to exercise judgment
		\end{enumerate}
	\end{frame}
	
	\begin{frame}{Q353: Explanation}
		\begin{block}{Correct Answer: C}
			\textbf{Why (C) is the best answer:}
			\begin{itemize}
				\item Usability drives adherence
				\item Handling rules must be clear
				\item Automation reduces reliance on users
			\end{itemize}
		\end{block}
		\begin{alertblock}{Critical Trap}
			A scheme that people cannot use is not a control.
		\end{alertblock}
	\end{frame}
	
	\begin{frame}{Q354: Legal Hold}
		\begin{block}{Question}
			\small A retention policy requires deletion of emails after 3 years, but the organization receives notice of litigation involving older emails. What should happen?
		\end{block}
		\begin{enumerate}[(A)]
			\footnotesize
			\item Continue deletion according to policy, since the retention schedule is the approved procedure and avoids storing excess data
			\item Move the relevant emails to the personal storage of the custodians, so that organizational systems still comply with the retention policy
			\item Delete the emails as scheduled but retain a summary document describing their content for use in the litigation
			\item Suspend deletion of relevant data under a legal hold managed by legal counsel
		\end{enumerate}
	\end{frame}
	
	\begin{frame}{Q354: Explanation}
		\begin{block}{Correct Answer: D}
			\textbf{Why (D) is the best answer:}
			\begin{itemize}
				\item Legal obligations override routine deletion
				\item Legal counsel directs the hold
				\item Destroying evidence creates liability
			\end{itemize}
		\end{block}
		\begin{alertblock}{Critical Trap}
			A legal hold overrides the retention schedule.
		\end{alertblock}
	\end{frame}
	
	\begin{frame}{Q355: Collecting Just in Case}
		\begin{block}{Question}
			\small A project proposes collecting additional customer attributes ``in case they are useful later.'' Which principle is MOST relevant?
		\end{block}
		\begin{enumerate}[(A)]
			\footnotesize
			\item Data minimization: collect only data necessary for a defined, legitimate purpose
			\item Data redundancy: retain multiple copies of each attribute so that future analytical needs can be met without recollection
			\item Data portability: allow customers to move their data so that additional attributes can be reused across services
			\item Data availability: keep all collected attributes online and accessible so that future projects can use them without delay
		\end{enumerate}
	\end{frame}
	
	\begin{frame}{Q355: Explanation}
		\begin{block}{Correct Answer: A}
			\textbf{Why (A) is the best answer:}
			\begin{itemize}
				\item Collect only what is needed
				\item Purpose must be defined
				\item Excess data is excess risk
			\end{itemize}
		\end{block}
		\begin{alertblock}{Critical Trap}
			Data you do not hold cannot be breached.
		\end{alertblock}
	\end{frame}
	
	\begin{frame}{Q356: Production Data in Test}
		\begin{block}{Question}
			\small Developers use copies of production data in test environments with weaker controls. What is the BEST mitigation?
		\end{block}
		\begin{enumerate}[(A)]
			\footnotesize
			\item Apply the same password policy to test environments as to production, which brings their authentication in line
			\item Allow the practice provided developers sign confidentiality agreements acknowledging the sensitivity of the data
			\item Use masked, tokenized, or synthetic data in non-production, allowing production data only under approved and controlled exceptions
			\item Restrict test environment access to one named developer, so that exposure of production data is limited to a single person
		\end{enumerate}
	\end{frame}
	
	\begin{frame}{Q356: Explanation}
		\begin{block}{Correct Answer: C}
			\textbf{Why (C) is the best answer:}
			\begin{itemize}
				\item Remove sensitivity from test data
				\item Exceptions require approval
				\item Controls should match data sensitivity
			\end{itemize}
		\end{block}
		\begin{alertblock}{Critical Trap}
			Sensitive data in weak environments is a breach waiting to happen.
		\end{alertblock}
	\end{frame}
	
	\begin{frame}{Q357: Ownerless Dataset}
		\begin{block}{Question}
			\small A critical dataset has no assigned owner, and each team assumes someone else is responsible. What is the MOST significant risk?
		\end{block}
		\begin{enumerate}[(A)]
			\footnotesize
			\item Storage costs will rise, because unowned datasets cannot be archived or deleted without an approver
			\item System performance will degrade, because unowned data is not tuned, indexed, or maintained
			\item No one is accountable for classification, access decisions, quality, or retention, so risk goes unmanaged
			\item Licensing fees will increase, because vendors charge more for datasets without a registered business owner
		\end{enumerate}
	\end{frame}
	
	\begin{frame}{Q357: Explanation}
		\begin{block}{Correct Answer: C}
			\textbf{Why (C) is the best answer:}
			\begin{itemize}
				\item Ownership drives decisions
				\item Access approvals need an owner
				\item Unowned data is unmanaged data
			\end{itemize}
		\end{block}
		\begin{alertblock}{Critical Trap}
			Unowned data is unmanaged risk.
		\end{alertblock}
	\end{frame}
	
	\begin{frame}{Q358: Conflicting Reports}
		\begin{block}{Question}
			\small Reports from a data warehouse regularly conflict with source system figures. Which approach BEST addresses this?
		\end{block}
		\begin{enumerate}[(A)]
			\footnotesize
			\item Implement data quality controls and lineage tracing with reconciliation and ownership at each stage
			\item Replace the business intelligence tool with a different product
			\item Instruct users to rely only on the source system
			\item Add more reports so users can cross-check them manually
		\end{enumerate}
	\end{frame}
	
	\begin{frame}{Q358: Explanation}
		\begin{block}{Correct Answer: A}
			\textbf{Why (A) is the best answer:}
			\begin{itemize}
				\item Trace data from source to report
				\item Reconcile at each stage
				\item Assign ownership of quality
			\end{itemize}
		\end{block}
		\begin{alertblock}{Critical Trap}
			Trust in data comes from lineage and reconciliation.
		\end{alertblock}
	\end{frame}
	
	\begin{frame}{Q359: Exiting a Cloud Provider}
		\begin{block}{Question}
			\small How can an organization BEST ensure data is irrecoverable when exiting a cloud provider whose physical media it cannot inspect?
		\end{block}
		\begin{enumerate}[(A)]
			\footnotesize
			\item Ask the provider to delete the files and rely on its confirmation email as evidence of deletion
			\item Overwrite the data once using the provider's management portal before closing the account
			\item Use cryptographic erasure by destroying keys the organization controls, supported by contractual deletion obligations and evidence
			\item Require physical destruction of the servers that stored the data, witnessed by the customer at the provider's data center
		\end{enumerate}
	\end{frame}
	
	\begin{frame}{Q359: Explanation}
		\begin{block}{Correct Answer: C}
			\textbf{Why (C) is the best answer:}
			\begin{itemize}
				\item Key destruction renders data unreadable
				\item Customer control of keys is essential
				\item Contracts require evidence of deletion
			\end{itemize}
		\end{block}
		\begin{alertblock}{Critical Trap}
			When you cannot destroy the media, destroy the keys.
		\end{alertblock}
	\end{frame}
	
	\begin{frame}{Q360: Cross-Border Processing}
		\begin{block}{Question}
			\small A business wants to use a global analytics service that processes personal data in another region. What should be assessed FIRST?
		\end{block}
		\begin{enumerate}[(A)]
			\footnotesize
			\item Performance benchmarks of the analytics service, to confirm it can process the expected data volumes
			\item Cost savings from regional pricing, to confirm the business case before any legal review begins
			\item The legal basis and transfer mechanisms for cross-border processing, together with the provider's safeguards
			\item Whether the data can be compressed before transfer, so that transmission cost and exposure are minimized
		\end{enumerate}
	\end{frame}
	
	\begin{frame}{Q360: Explanation}
		\begin{block}{Correct Answer: C}
			\textbf{Why (C) is the best answer:}
			\begin{itemize}
				\item Privacy law governs where data may go
				\item Transfer mechanisms must exist
				\item Evaluate provider safeguards
			\end{itemize}
		\end{block}
		\begin{alertblock}{Critical Trap}
			Legality comes before performance.
		\end{alertblock}
	\end{frame}
	
	\begin{frame}{Q361: DLP Bypass}
		\begin{block}{Question}
			\small DLP generates many false positives, and users bypass it by using personal email. What is the BEST response?
		\end{block}
		\begin{enumerate}[(A)]
			\footnotesize
			\item Tune policies with business input, combine with user education, and address the workflow needs that cause users to seek alternatives
			\item Disable DLP to improve productivity until the false-positive rate falls to an acceptable level
			\item Increase blocking sensitivity so that fewer sensitive files can leave through any channel, including personal email
			\item Replace DLP with a stronger perimeter firewall that blocks all external email destinations by default
		\end{enumerate}
	\end{frame}
	
	\begin{frame}{Q361: Explanation}
		\begin{block}{Correct Answer: A}
			\textbf{Why (A) is the best answer:}
			\begin{itemize}
				\item Tune policies with business input
				\item Educate users
				\item Fix the workflow need
			\end{itemize}
		\end{block}
		\begin{alertblock}{Critical Trap}
			Controls that block work get bypassed.
		\end{alertblock}
	\end{frame}
	
	\begin{frame}{Q362: Tokenization Versus Encryption}
		\begin{block}{Question}
			\small Which statement BEST differentiates tokenization from encryption?
		\end{block}
		\begin{enumerate}[(A)]
			\footnotesize
			\item Tokenization is reversible using an encryption key held by the application, whereas encryption cannot be reversed once applied
			\item Tokenization replaces data with surrogate values mapped in a secure vault, while encryption mathematically transforms data and is reversible with a key
			\item Both are hashing techniques; they differ only in output format and are interchangeable for protecting regulated data
			\item Tokenization always requires more computing power than encryption, which is why it is used only for small datasets
		\end{enumerate}
	\end{frame}
	
	\begin{frame}{Q362: Explanation}
		\begin{block}{Correct Answer: B}
			\textbf{Why (B) is the best answer:}
			\begin{itemize}
				\item Tokens have no mathematical relationship to data
				\item Vault mapping enables detokenization
				\item Encryption depends on key protection
			\end{itemize}
		\end{block}
		\begin{alertblock}{Critical Trap}
			Tokens are substitutes; ciphertext is transformed data.
		\end{alertblock}
	\end{frame}
	
	\begin{frame}{Q363: Erasure and Backups}
		\begin{block}{Question}
			\small A data subject requests erasure, but the data also exists in backups that cannot be selectively modified. What is the BEST approach?
		\end{block}
		\begin{enumerate}[(A)]
			\footnotesize
			\item Refuse the request, because backups are always exempt from erasure rights under privacy regulations
			\item Delete all backup sets immediately, so that complete erasure is guaranteed regardless of recovery consequences
			\item Delete from live systems, document the backup limitation, ensure the data is not restored or used, and let it expire under the backup retention schedule
			\item Ignore the request until a backup is next restored, then delete the data from the restored system
		\end{enumerate}
	\end{frame}
	
	\begin{frame}{Q363: Explanation}
		\begin{block}{Correct Answer: C}
			\textbf{Why (C) is the best answer:}
			\begin{itemize}
				\item Erase from active systems first
				\item Document technical limits
				\item Prevent use of backup copies
			\end{itemize}
		\end{block}
		\begin{alertblock}{Critical Trap}
			Handle backup limits transparently, not by ignoring them.
		\end{alertblock}
	\end{frame}
	
	\begin{frame}{Q364: Shadow Data}
		\begin{block}{Question}
			\small Discovery tools find sensitive data in unmanaged file shares and personal cloud storage. What is the BEST action?
		\end{block}
		\begin{enumerate}[(A)]
			\footnotesize
			\item Delete all discovered data immediately, since unmanaged data has no business justification
			\item Block all access to cloud storage services across the organization until a full inventory is completed
			\item Accept the risk and issue a reminder, since users are responsible for the storage locations they choose
			\item Inventory and classify the data, determine ownership, move or protect it under approved controls, and address the behaviors that created it
		\end{enumerate}
	\end{frame}
	
	\begin{frame}{Q364: Explanation}
		\begin{block}{Correct Answer: D}
			\textbf{Why (D) is the best answer:}
			\begin{itemize}
				\item Discover, classify, and assign owners
				\item Move data into controlled locations
				\item Fix the root behavior
			\end{itemize}
		\end{block}
		\begin{alertblock}{Critical Trap}
			Discovery is the start, not the end.
		\end{alertblock}
	\end{frame}
	
	\begin{frame}{Q365: Why Map Data}
		\begin{block}{Question}
			\small Why is a data inventory and map foundational to data lifecycle risk management?
		\end{block}
		\begin{enumerate}[(A)]
			\footnotesize
			\item It replaces the need for classification, because location and flow information already indicate sensitivity
			\item It is required only for backup planning, since other controls rely on classification rather than location
			\item Controls, retention, privacy obligations, and incident response all depend on knowing where data resides and how it flows
			\item It reduces the need for periodic access reviews, because the flows already show who uses each dataset
		\end{enumerate}
	\end{frame}
	
	\begin{frame}{Q365: Explanation}
		\begin{block}{Correct Answer: C}
			\textbf{Why (C) is the best answer:}
			\begin{itemize}
				\item You cannot protect unknown data
				\item Privacy obligations depend on flows
				\item Incident response needs data location
			\end{itemize}
		\end{block}
		\begin{alertblock}{Critical Trap}
			You cannot manage data you cannot locate.
		\end{alertblock}
	\end{frame}
	
	\begin{frame}{Q366: Sub-Processors}
		\begin{block}{Question}
			\small A processor sub-contracts part of its service to another vendor that handles your customer data. Which control is MOST important?
		\end{block}
		\begin{enumerate}[(A)]
			\footnotesize
			\item No additional action is needed, because the primary vendor remains fully liable for its subcontractors
			\item Ask the sub-vendor to sign a confidentiality agreement, which is sufficient assurance of protection
			\item Require disclosure of the vendor's pricing and margin structure so that cost transparency is maintained across the chain
			\item Contractual requirements for sub-processor approval and flow-down of security and privacy obligations, with the ability to assess compliance
		\end{enumerate}
	\end{frame}
	
	\begin{frame}{Q366: Explanation}
		\begin{block}{Correct Answer: D}
			\textbf{Why (D) is the best answer:}
			\begin{itemize}
				\item Obligations must flow down
				\item Approval rights over sub-processors
				\item Right to assess compliance
			\end{itemize}
		\end{block}
		\begin{alertblock}{Critical Trap}
			Your obligations extend down the supply chain.
		\end{alertblock}
	\end{frame}
	
	\begin{frame}{Q367: Inconsistent Retention}
		\begin{block}{Question}
			\small Different departments keep data for varying periods with no consistent basis. What is the BEST approach?
		\end{block}
		\begin{enumerate}[(A)]
			\footnotesize
			\item Keep all data indefinitely to avoid any risk of noncompliance with retention obligations
			\item Allow each data owner to choose a retention period based on the needs of their own department
			\item Delete all data after one year to minimize exposure, regardless of legal or business requirements
			\item Establish a retention schedule based on legal, regulatory, and business requirements with defined disposal procedures
		\end{enumerate}
	\end{frame}
	
	\begin{frame}{Q367: Explanation}
		\begin{block}{Correct Answer: D}
			\textbf{Why (D) is the best answer:}
			\begin{itemize}
				\item Retention must reflect obligations
				\item Schedule and disposal are linked
				\item Consistency supports compliance
			\end{itemize}
		\end{block}
		\begin{alertblock}{Critical Trap}
			Retention is driven by requirements, not preference.
		\end{alertblock}
	\end{frame}
	
	\begin{frame}{Q368: Long-Term Archives}
		\begin{block}{Question}
			\small Archived records must remain readable and trustworthy for 15 years. Which risk is OFTEN overlooked?
		\end{block}
		\begin{enumerate}[(A)]
			\footnotesize
			\item Higher storage cost per gigabyte over time, which can make long-term archiving unaffordable
			\item Backup window limitations, because long-term archives require additional nightly capacity
			\item Network bandwidth constraints, which limit the speed at which archived records can be retrieved
			\item Format and media obsolescence and loss of integrity or key access over time
		\end{enumerate}
	\end{frame}
	
	\begin{frame}{Q368: Explanation}
		\begin{block}{Correct Answer: D}
			\textbf{Why (D) is the best answer:}
			\begin{itemize}
				\item Formats and media become obsolete
				\item Keys must remain available
				\item Integrity must be verifiable
			\end{itemize}
		\end{block}
		\begin{alertblock}{Critical Trap}
			Archives fail through obsolescence and lost keys.
		\end{alertblock}
	\end{frame}
		
	% ============================================================
	\section{5. SDLC and Emerging Technologies}
	% ============================================================
	
	\begin{frame}{Q369: Security Requirements Missing}
		\begin{block}{Question}
			\small A project reaches testing and no security requirements were defined. What is the PRIMARY consequence?
		\end{block}
		\begin{enumerate}[(A)]
			\footnotesize
			\item Testers will have fewer test cases to execute, which will shorten the testing phase and the schedule
			\item The project will automatically fail its next audit, because missing requirements are always an audit finding
			\item Security gaps will be found late, when remediation is costlier and may be accepted under schedule pressure
			\item Developers will be unable to use security tools, because those tools are configured from the requirements document
		\end{enumerate}
	\end{frame}
	
	\begin{frame}{Q369: Explanation}
		\begin{block}{Correct Answer: C}
			\textbf{Why (C) is the best answer:}
			\begin{itemize}
				\item Define requirements early
				\item Late fixes cost more
				\item Schedule pressure leads to acceptance
			\end{itemize}
		\end{block}
		\begin{alertblock}{Critical Trap}
			Security requirements belong at the start.
		\end{alertblock}
	\end{frame}
	
	\begin{frame}{Q370: SAST and DAST}
		\begin{block}{Question}
			\small Which statement BEST describes how SAST and DAST complement each other?
		\end{block}
		\begin{enumerate}[(A)]
			\footnotesize
			\item SAST analyzes source code for flaws early, while DAST tests the running application for exploitable behavior, giving broader coverage together
			\item SAST tests running applications and DAST analyzes source code, so both are needed to cover the full pipeline
			\item They are redundant because each finds the same flaws, so one of them is sufficient for assurance
			\item Both evaluate only infrastructure configuration and complement each other by covering different network layers
		\end{enumerate}
	\end{frame}
	
	\begin{frame}{Q370: Explanation}
		\begin{block}{Correct Answer: A}
			\textbf{Why (A) is the best answer:}
			\begin{itemize}
				\item SAST reviews code
				\item DAST tests the running application
				\item Together they cover more defects
			\end{itemize}
		\end{block}
		\begin{alertblock}{Critical Trap}
			Static and dynamic testing find different flaws.
		\end{alertblock}
	\end{frame}
	
	\begin{frame}{Q371: Unknown Library Exposure}
		\begin{block}{Question}
			\small A critical vulnerability in an open-source library is announced. The organization cannot quickly tell which applications use it. What capability is MISSING?
		\end{block}
		\begin{enumerate}[(A)]
			\footnotesize
			\item A software bill of materials and dependency inventory to identify affected components
			\item A larger security operations center, so that more analysts can search for the library across applications
			\item Stronger network segmentation, so that vulnerable libraries cannot be reached from untrusted networks
			\item Annual penetration testing, which would have identified the library as a vulnerable component
		\end{enumerate}
	\end{frame}
	
	\begin{frame}{Q371: Explanation}
		\begin{block}{Correct Answer: A}
			\textbf{Why (A) is the best answer:}
			\begin{itemize}
				\item You must know your components
				\item SBOM speeds impact assessment
				\item Supply chain visibility is a control
			\end{itemize}
		\end{block}
		\begin{alertblock}{Critical Trap}
			If you cannot list your components, you cannot patch them.
		\end{alertblock}
	\end{frame}
	
	\begin{frame}{Q372: Keys in Repository}
		\begin{block}{Question}
			\small A scan finds API keys committed to a code repository. What is the BEST response?
		\end{block}
		\begin{enumerate}[(A)]
			\footnotesize
			\item Delete the commit from repository history, since the keys are then no longer accessible to others
			\item Make the repository private, which restricts access to the keys to authorized staff
			\item Warn developers by email not to repeat the practice and add the rule to the coding standard
			\item Revoke and rotate the exposed credentials, then implement secrets management and pre-commit scanning
		\end{enumerate}
	\end{frame}
	
	\begin{frame}{Q372: Explanation}
		\begin{block}{Correct Answer: D}
			\textbf{Why (D) is the best answer:}
			\begin{itemize}
				\item Exposed credentials must be rotated
				\item Use a secrets manager
				\item Prevent recurrence with scanning
			\end{itemize}
		\end{block}
		\begin{alertblock}{Critical Trap}
			Once exposed, a secret must be rotated.
		\end{alertblock}
	\end{frame}
	
	\begin{frame}{Q373: Deferred Security in Agile}
		\begin{block}{Question}
			\small In an agile project, security reviews and documentation are repeatedly deferred. What is the BEST way to address this?
		\end{block}
		\begin{enumerate}[(A)]
			\footnotesize
			\item Move the project to a waterfall approach, where security phases are fixed and cannot be skipped
			\item Integrate security activities into sprints through the definition of done, automated checks, and risk-based acceptance criteria
			\item Add a single security review at the end of the project to ensure that all issues are found before release
			\item Remove documentation requirements entirely, to keep the team focused on delivering working software
		\end{enumerate}
	\end{frame}
	
	\begin{frame}{Q373: Explanation}
		\begin{block}{Correct Answer: B}
			\textbf{Why (B) is the best answer:}
			\begin{itemize}
				\item Make security part of done
				\item Automate checks in the pipeline
				\item Use risk-based criteria
			\end{itemize}
		\end{block}
		\begin{alertblock}{Critical Trap}
			Agile needs security built in, not bolted on.
		\end{alertblock}
	\end{frame}
	
	\begin{frame}{Q374: Self-Approved Code}
		\begin{block}{Question}
			\small The developer who wrote a change also approves their own pull request. What control gap exists?
		\end{block}
		\begin{enumerate}[(A)]
			\footnotesize
			\item Insufficient unit test coverage, which allows defects to pass unnoticed through the pipeline
			\item Inconsistent code style across the repository, which makes automated review less reliable
			\item Lack of independent review, increasing the risk of errors or malicious code reaching production
			\item Repository storage limits, which prevent reviewers from keeping a history of approvals
		\end{enumerate}
	\end{frame}
	
	\begin{frame}{Q374: Explanation}
		\begin{block}{Correct Answer: C}
			\textbf{Why (C) is the best answer:}
			\begin{itemize}
				\item Independent review catches errors
				\item Self-approval defeats segregation
				\item Enforce reviewer rules in tooling
			\end{itemize}
		\end{block}
		\begin{alertblock}{Critical Trap}
			Nobody should approve their own work.
		\end{alertblock}
	\end{frame}
	
	\begin{frame}{Q375: Developer-Run UAT}
		\begin{block}{Question}
			\small User acceptance testing is performed by the same team that developed the application. What is the MAIN concern?
		\end{block}
		\begin{enumerate}[(A)]
			\footnotesize
			\item Developers will complete testing too slowly, because they must switch between development and testing roles
			\item Lack of independence and business perspective reduces the likelihood of finding requirement or usability gaps
			\item UAT must always use production data, which developers are not authorized to access
			\item Documentation will be incomplete, because developers do not write test scripts for their own code
		\end{enumerate}
	\end{frame}
	
	\begin{frame}{Q375: Explanation}
		\begin{block}{Correct Answer: B}
			\textbf{Why (B) is the best answer:}
			\begin{itemize}
				\item UAT validates business need
				\item Independence adds objectivity
				\item Business users should execute UAT
			\end{itemize}
		\end{block}
		\begin{alertblock}{Critical Trap}
			UAT belongs to the business.
		\end{alertblock}
	\end{frame}
	
	\begin{frame}{Q376: Machine Learning Credit Model}
		\begin{block}{Question}
			\small A bank deploys a machine learning credit decision model. Which governance element is MOST critical?
		\end{block}
		\begin{enumerate}[(A)]
			\footnotesize
			\item Training on the largest available dataset, since broader data always improves predictive performance and fairness
			\item Model validation, ongoing monitoring for drift and bias, and clear accountability and explainability appropriate to the decisions
			\item The reputation of the model vendor, since an established vendor's validation is sufficient assurance
			\item Measuring accuracy once at launch, since a model that performs well initially will continue to do so
		\end{enumerate}
	\end{frame}
	
	\begin{frame}{Q376: Explanation}
		\begin{block}{Correct Answer: B}
			\textbf{Why (B) is the best answer:}
			\begin{itemize}
				\item Validate before and after deployment
				\item Monitor for drift and bias
				\item Assign accountability
			\end{itemize}
		\end{block}
		\begin{alertblock}{Critical Trap}
			Models need governance across their whole life.
		\end{alertblock}
	\end{frame}
	
	\begin{frame}{Q377: Training Data Manipulation}
		\begin{block}{Question}
			\small Which risk is SPECIFIC to machine learning systems trained on external data?
		\end{block}
		\begin{enumerate}[(A)]
			\footnotesize
			\item Unauthorized access to the training environment through weak administrator passwords, a risk that applies equally to any system
			\item Data poisoning, where manipulated training data degrades or biases model behavior
			\item Injection of malicious SQL statements into the application database, a risk common to any web application
			\item Memory corruption flaws in the operating system hosting the model, which exist regardless of the type of model
		\end{enumerate}
	\end{frame}
	
	\begin{frame}{Q377: Explanation}
		\begin{block}{Correct Answer: B}
			\textbf{Why (B) is the best answer:}
			\begin{itemize}
				\item Training data shapes behavior
				\item Attackers can corrupt inputs
				\item Validate data sources
			\end{itemize}
		\end{block}
		\begin{alertblock}{Critical Trap}
			Model integrity depends on data integrity.
		\end{alertblock}
	\end{frame}
	
	\begin{frame}{Q378: Generative AI Leakage}
		\begin{block}{Question}
			\small Employees paste confidential code and customer data into a public generative AI tool. What is the BEST combination of controls?
		\end{block}
		\begin{enumerate}[(A)]
			\footnotesize
			\item Ban all generative AI tools through a policy statement, relying on employees to comply
			\item Trust employees' judgment about what to share, supported by an annual reminder of confidentiality obligations
			\item Rely on the AI vendor's terms of service, which state that submitted data is not used for training
			\item Establish an acceptable use policy, provide approved AI tools with data protections, and use technical controls such as DLP and access restrictions
		\end{enumerate}
	\end{frame}
	
	\begin{frame}{Q378: Explanation}
		\begin{block}{Correct Answer: D}
			\textbf{Why (D) is the best answer:}
			\begin{itemize}
				\item Policy needs enforcement
				\item Offer a safe approved alternative
				\item Use DLP to stop leakage
			\end{itemize}
		\end{block}
		\begin{alertblock}{Critical Trap}
			Provide a safe path and enforce it.
		\end{alertblock}
	\end{frame}
	
	\begin{frame}{Q379: Default Credentials on IoT}
		\begin{block}{Question}
			\small A manufacturer deploys thousands of IoT sensors with default credentials and no update mechanism. What is the MOST significant risk?
		\end{block}
		\begin{enumerate}[(A)]
			\footnotesize
			\item Unmanaged, easily compromised devices that provide network entry points and cannot be patched
			\item Increased electricity consumption across the plant, which raises operating costs
			\item Data format incompatibility with analytics tools, which prevents collected data from being used
			\item Gradual drift in sensor accuracy, which reduces the quality of production data over time
		\end{enumerate}
	\end{frame}
	
	\begin{frame}{Q379: Explanation}
		\begin{block}{Correct Answer: A}
			\textbf{Why (A) is the best answer:}
			\begin{itemize}
				\item Default credentials are well known
				\item Unpatchable devices stay vulnerable
				\item Segment and inventory devices
			\end{itemize}
		\end{block}
		\begin{alertblock}{Critical Trap}
			Unmanageable devices are permanent entry points.
		\end{alertblock}
	\end{frame}
	
	\begin{frame}{Q380: Smart Contract Deployment}
		\begin{block}{Question}
			\small Before deploying a smart contract that manages financial assets, which activity is MOST important?
		\end{block}
		\begin{enumerate}[(A)]
			\footnotesize
			\item A marketing campaign to attract users, since adoption determines the contract's success
			\item Independent code audit and extensive testing, since deployed code may be immutable and flaws irreversible
			\item Raising transaction fees so that the contract executes with priority on the network
			\item Selecting a faster block time, so that transactions confirm sooner and fraud is detected earlier
		\end{enumerate}
	\end{frame}
	
	\begin{frame}{Q380: Explanation}
		\begin{block}{Correct Answer: B}
			\textbf{Why (B) is the best answer:}
			\begin{itemize}
				\item Deployed code is hard to change
				\item Audit before release
				\item Test edge cases thoroughly
			\end{itemize}
		\end{block}
		\begin{alertblock}{Critical Trap}
			Immutable code demands pre-release assurance.
		\end{alertblock}
	\end{frame}
	
	\begin{frame}{Q381: Low-Code Sprawl}
		\begin{block}{Question}
			\small Business users build workflows with low-code tools that connect to sensitive data without IT involvement. What is the BEST governance response?
		\end{block}
		\begin{enumerate}[(A)]
			\footnotesize
			\item Provide a governed platform with guardrails, an inventory, data access controls, and review for higher-risk applications
			\item Prohibit all low-code development and require requests to go through the standard IT project process
			\item Allow users to self-certify the security of their applications through an annual questionnaire
			\item Require IT to build every workflow, so that all applications follow the standard development lifecycle
		\end{enumerate}
	\end{frame}
	
	\begin{frame}{Q381: Explanation}
		\begin{block}{Correct Answer: A}
			\textbf{Why (A) is the best answer:}
			\begin{itemize}
				\item Offer a governed platform
				\item Maintain an inventory
				\item Review by risk level
			\end{itemize}
		\end{block}
		\begin{alertblock}{Critical Trap}
			Enable citizen developers safely.
		\end{alertblock}
	\end{frame}
	
	\begin{frame}{Q382: Bot Credentials}
		\begin{block}{Question}
			\small Robotic process automation bots log into multiple systems using a shared admin account. What is the BEST improvement?
		\end{block}
		\begin{enumerate}[(A)]
			\footnotesize
			\item Increase the length and complexity of the shared password and rotate it quarterly
			\item Assign unique, least-privilege bot identities with credentials in a vault, logging, and a named owner
			\item Run bots only during overnight windows when fewer users are active, to reduce exposure
			\item Remove bot logging to improve performance, since bot actions are already predictable
		\end{enumerate}
	\end{frame}
	
	\begin{frame}{Q382: Explanation}
		\begin{block}{Correct Answer: B}
			\textbf{Why (B) is the best answer:}
			\begin{itemize}
				\item Bots need unique identities
				\item Vault the credentials
				\item Log bot actions
			\end{itemize}
		\end{block}
		\begin{alertblock}{Critical Trap}
			Bots are identities and need least privilege.
		\end{alertblock}
	\end{frame}
	
	\begin{frame}{Q383: Unscanned Container Images}
		\begin{block}{Question}
			\small Containers are built from public base images with no scanning. What is the PRIMARY risk?
		\end{block}
		\begin{enumerate}[(A)]
			\footnotesize
			\item Slower container startup times, because public images are larger than curated internal images
			\item Higher licensing fees, because public images carry commercial license obligations
			\item Increased network traffic, because images must be pulled from external registries on every deployment
			\item Inherited vulnerabilities or malicious components in public base images are deployed at scale across every container built from them
		\end{enumerate}
	\end{frame}
	
	\begin{frame}{Q383: Explanation}
		\begin{block}{Correct Answer: D}
			\textbf{Why (D) is the best answer:}
			\begin{itemize}
				\item Base images can be vulnerable
				\item Scan images in the pipeline
				\item Use trusted registries
			\end{itemize}
		\end{block}
		\begin{alertblock}{Critical Trap}
			Scan what you build from.
		\end{alertblock}
	\end{frame}
	
	\begin{frame}{Q384: Exposed APIs}
		\begin{block}{Question}
			\small Which is the MOST important control for externally exposed APIs?
		\end{block}
		\begin{enumerate}[(A)]
			\footnotesize
			\item Strong authentication and authorization on each request, with input validation and rate limiting
			\item Obscure URL paths that are hard to guess, so that only intended consumers discover the API
			\item Keeping the API documentation hidden from the public, so that attackers cannot learn how the API works
			\item Running the API on a non-standard port, so that automated scanners do not detect it
		\end{enumerate}
	\end{frame}
	
	\begin{frame}{Q384: Explanation}
		\begin{block}{Correct Answer: A}
			\textbf{Why (A) is the best answer:}
			\begin{itemize}
				\item Authenticate and authorize every call
				\item Validate input
				\item Rate limit to resist abuse
			\end{itemize}
		\end{block}
		\begin{alertblock}{Critical Trap}
			Obscurity is not authorization.
		\end{alertblock}
	\end{frame}
	
	\begin{frame}{Q385: Crypto Agility}
		\begin{block}{Question}
			\small Why should organizations consider cryptographic agility in light of quantum computing?
		\end{block}
		\begin{enumerate}[(A)]
			\footnotesize
			\item Because quantum computing makes hashing impossible, so all integrity controls must be replaced immediately
			\item To replace vulnerable algorithms without major redesign, since long-lived confidential data may be harvested now and decrypted later
			\item Because quantum computers will make symmetric encryption trivially breakable, so all data must be re-encrypted this year
			\item Because all current systems are already quantum-safe, so only governance of algorithm choices is needed
		\end{enumerate}
	\end{frame}
	
	\begin{frame}{Q385: Explanation}
		\begin{block}{Correct Answer: B}
			\textbf{Why (B) is the best answer:}
			\begin{itemize}
				\item Algorithms may need replacement
				\item Harvest now, decrypt later is a risk
				\item Design for replaceable crypto
			\end{itemize}
		\end{block}
		\begin{alertblock}{Critical Trap}
			Plan for algorithm change before it is forced.
		\end{alertblock}
	\end{frame}
		
	% ============================================================
	\section{6. Information Security Principles}
	% ============================================================
	
	\begin{frame}{Q386: Access Accumulation}
		\begin{block}{Question}
			\small An employee has transferred departments three times and retains access from every previous role. Which control would have BEST prevented this?
		\end{block}
		\begin{enumerate}[(A)]
			\footnotesize
			\item Access review triggered by role change, plus periodic recertification by business owners
			\item Stronger password complexity requirements, which would reduce the chance that old accounts are misused
			\item Annual security awareness training, which would remind employees to surrender access they no longer need
			\item Single sign-on across applications, which would make it easier to remove access from one place
		\end{enumerate}
	\end{frame}
	
	\begin{frame}{Q386: Explanation}
		\begin{block}{Correct Answer: A}
			\textbf{Why (A) is the best answer:}
			\begin{itemize}
				\item Role change must trigger review
				\item Owners recertify access
				\item Prevents privilege accumulation
			\end{itemize}
		\end{block}
		\begin{alertblock}{Critical Trap}
			Access must follow the role, not the history.
		\end{alertblock}
	\end{frame}
	
	\begin{frame}{Q387: Privileged Access}
		\begin{block}{Question}
			\small Which privileged access management capability BEST reduces risk from compromised administrator credentials?
		\end{block}
		\begin{enumerate}[(A)]
			\footnotesize
			\item Longer and more complex passwords for administrator accounts, changed at each employee departure
			\item Shared administrator accounts with scheduled rotation of passwords every thirty days
			\item Just-in-time privileged access with credential vaulting and session recording
			\item Email notification to the security team of every administrator login, reviewed weekly
		\end{enumerate}
	\end{frame}
	
	\begin{frame}{Q387: Explanation}
		\begin{block}{Correct Answer: C}
			\textbf{Why (C) is the best answer:}
			\begin{itemize}
				\item Limit standing privileges
				\item Vault credentials
				\item Record sessions for accountability
			\end{itemize}
		\end{block}
		\begin{alertblock}{Critical Trap}
			Reduce standing privilege and record what is done.
		\end{alertblock}
	\end{frame}
	
	\begin{frame}{Q388: Push Bombing}
		\begin{block}{Question}
			\small Attackers bombard users with push approval requests until one is accepted. What is the BEST mitigation?
		\end{block}
		\begin{enumerate}[(A)]
			\footnotesize
			\item Disable push-based MFA and rely on strong passwords, since fatigue attacks target only the approval step
			\item Increase password length requirements so that stolen credentials are insufficient on their own
			\item Switch from push approvals to SMS one-time codes, which are not susceptible to repeated prompts
			\item Phishing-resistant MFA or number matching, with rate limiting and user reporting
		\end{enumerate}
	\end{frame}
	
	\begin{frame}{Q388: Explanation}
		\begin{block}{Correct Answer: D}
			\textbf{Why (D) is the best answer:}
			\begin{itemize}
				\item Prompt fatigue is an attack method
				\item Use number matching or FIDO
				\item Rate limit and report
			\end{itemize}
		\end{block}
		\begin{alertblock}{Critical Trap}
			MFA must resist social engineering.
		\end{alertblock}
	\end{frame}
	
	\begin{frame}{Q389: Zero Trust Principle}
		\begin{block}{Question}
			\small Which statement BEST reflects zero trust implementation?
		\end{block}
		\begin{enumerate}[(A)]
			\footnotesize
			\item Users are trusted after VPN login because they are on the internal network
			\item A next-generation firewall at the perimeter provides sufficient protection
			\item MFA is required once per year for all employees
			\item Access is decided per request using identity, device posture, and context, with least privilege and continuous verification
		\end{enumerate}
	\end{frame}
	
	\begin{frame}{Q389: Explanation}
		\begin{block}{Correct Answer: D}
			\textbf{Why (D) is the best answer:}
			\begin{itemize}
				\item Never trust by location
				\item Verify identity, device, and context
				\item Apply least privilege continuously
			\end{itemize}
		\end{block}
		\begin{alertblock}{Critical Trap}
			Location is not trust.
		\end{alertblock}
	\end{frame}
	
	\begin{frame}{Q390: Keys Beside Data}
		\begin{block}{Question}
			\small Encryption keys are stored on the same server as the encrypted data. What is the MAIN issue?
		\end{block}
		\begin{enumerate}[(A)]
			\footnotesize
			\item System performance will degrade, because decryption on the same server consumes additional processing capacity
			\item Compromise of the server exposes both data and keys, defeating the protection; keys should be separately protected, for example in an HSM or KMS
			\item Keys will expire sooner when stored on production servers, causing frequent encryption failures
			\item Regulations prohibit encryption on application servers, so data must be encrypted in the database tier only
		\end{enumerate}
	\end{frame}
	
	\begin{frame}{Q390: Explanation}
		\begin{block}{Correct Answer: B}
			\textbf{Why (B) is the best answer:}
			\begin{itemize}
				\item Keys must be separate from data
				\item Use HSM or KMS
				\item Control key access tightly
			\end{itemize}
		\end{block}
		\begin{alertblock}{Critical Trap}
			Encryption is only as strong as key protection.
		\end{alertblock}
	\end{frame}
	
	\begin{frame}{Q391: Defense in Depth}
		\begin{block}{Question}
			\small Which design BEST demonstrates defense in depth for a sensitive database?
		\end{block}
		\begin{enumerate}[(A)]
			\footnotesize
			\item One very strong perimeter firewall with advanced threat detection, which blocks attacks before they reach the database
			\item Multiple identical firewalls in series, so that a failure in one is covered by the next
			\item Layered controls including segmentation, access control, encryption, monitoring, and backups
			\item Strong complex passwords on all database accounts, changed regularly and stored in a vault
		\end{enumerate}
	\end{frame}
	
	\begin{frame}{Q391: Explanation}
		\begin{block}{Correct Answer: C}
			\textbf{Why (C) is the best answer:}
			\begin{itemize}
				\item Multiple layers of control
				\item Different control types
				\item One failure does not expose the asset
			\end{itemize}
		\end{block}
		\begin{alertblock}{Critical Trap}
			No single control should stand alone.
		\end{alertblock}
	\end{frame}
	
	\begin{frame}{Q392: Premature Eradication}
		\begin{block}{Question}
			\small During an active intrusion, the team wants to immediately wipe and rebuild affected servers. What is the concern?
		\end{block}
		\begin{enumerate}[(A)]
			\footnotesize
			\item Rebuilding is always preferred to investigation, since it removes the attacker immediately and limits further damage
			\item Premature eradication may destroy forensic evidence and miss attacker persistence if scope is not understood; containment and investigation come first
			\item Wiping servers violates cloud service terms, which prohibit customers from destroying provider-managed resources
			\item Servers cannot be rebuilt during business hours, so the activity must wait for the next maintenance window
		\end{enumerate}
	\end{frame}
	
	\begin{frame}{Q392: Explanation}
		\begin{block}{Correct Answer: B}
			\textbf{Why (B) is the best answer:}
			\begin{itemize}
				\item Contain first
				\item Preserve evidence
				\item Understand scope before eradicating
			\end{itemize}
		\end{block}
		\begin{alertblock}{Critical Trap}
			Investigate scope before you eradicate.
		\end{alertblock}
	\end{frame}
	
	\begin{frame}{Q393: Vulnerability Prioritization}
		\begin{block}{Question}
			\small Thousands of vulnerabilities exist. What is the BEST prioritization approach?
		\end{block}
		\begin{enumerate}[(A)]
			\footnotesize
			\item Prioritize strictly by CVSS base score, since it is the standardized and objective measure of severity
			\item Risk-based prioritization using exploitability, asset criticality, exposure, and compensating controls, not CVSS alone
			\item Remediate the oldest vulnerabilities first, since they have been exposed to attackers for the longest time
			\item Remediate the easiest vulnerabilities first, so that the largest number of findings are closed quickly
		\end{enumerate}
	\end{frame}
	
	\begin{frame}{Q393: Explanation}
		\begin{block}{Correct Answer: B}
			\textbf{Why (B) is the best answer:}
			\begin{itemize}
				\item Combine exploitability and impact
				\item Consider asset criticality
				\item Account for compensating controls
			\end{itemize}
		\end{block}
		\begin{alertblock}{Critical Trap}
			Severity scores need business context.
		\end{alertblock}
	\end{frame}
	
	\begin{frame}{Q394: Awareness Effectiveness}
		\begin{block}{Question}
			\small Which metric BEST shows the effectiveness of security awareness training?
		\end{block}
		\begin{enumerate}[(A)]
			\footnotesize
			\item Behavior metrics such as phishing report rates and declining click rates over time
			\item Training completion percentage, which shows that the workforce has been reached by the program
			\item Total hours of training delivered, which indicates the depth of the awareness program
			\item Average quiz scores after the course, which show that employees have retained the material
		\end{enumerate}
	\end{frame}
	
	\begin{frame}{Q394: Explanation}
		\begin{block}{Correct Answer: A}
			\textbf{Why (A) is the best answer:}
			\begin{itemize}
				\item Measure behavior change
				\item Reporting indicates engagement
				\item Completion is not effectiveness
			\end{itemize}
		\end{block}
		\begin{alertblock}{Critical Trap}
			Measure behavior, not attendance.
		\end{alertblock}
	\end{frame}
	
	\begin{frame}{Q395: SIEM Without Use Cases}
		\begin{block}{Question}
			\small A SIEM collects logs from many sources, but no detection use cases are defined. What is the MAIN weakness?
		\end{block}
		\begin{enumerate}[(A)]
			\footnotesize
			\item Storage costs will exceed budget, because logs from all sources are retained without selection
			\item Log formats will be inconsistent, which prevents correlation across sources
			\item Analysts will need additional certification, because the platform cannot be used without vendor-specific training
			\item Collected data is not turned into detection of relevant threats; use cases aligned with risk scenarios are needed
		\end{enumerate}
	\end{frame}
	
	\begin{frame}{Q395: Explanation}
		\begin{block}{Correct Answer: D}
			\textbf{Why (D) is the best answer:}
			\begin{itemize}
				\item Logs alone do not detect
				\item Align use cases to risk scenarios
				\item Tune continuously
			\end{itemize}
		\end{block}
		\begin{alertblock}{Critical Trap}
			A SIEM without use cases is an expensive archive.
		\end{alertblock}
	\end{frame}
	
	\begin{frame}{Q396: Exposed Cloud Bucket}
		\begin{block}{Question}
			\small A company's IaaS storage bucket is publicly exposed due to a misconfiguration. Who is MOST likely responsible?
		\end{block}
		\begin{enumerate}[(A)]
			\footnotesize
			\item The cloud provider in every case, since it operates the platform that hosts the storage service
			\item Both parties equally, regardless of the service model, since responsibility is always split evenly under the shared responsibility model
			\item The regulator that certified the provider, since certification implies that default configurations are secure
			\item The customer, since configuring access controls on its own data is the customer's responsibility under shared responsibility
		\end{enumerate}
	\end{frame}
	
	\begin{frame}{Q396: Explanation}
		\begin{block}{Correct Answer: D}
			\textbf{Why (D) is the best answer:}
			\begin{itemize}
				\item Customers configure their own access
				\item Providers secure the infrastructure
				\item Misconfiguration is a leading cloud risk
			\end{itemize}
		\end{block}
		\begin{alertblock}{Critical Trap}
			Shared responsibility does not remove customer duties.
		\end{alertblock}
	\end{frame}
	
	\begin{frame}{Q397: Insider Threat}
		\begin{block}{Question}
			\small Which combination BEST mitigates insider threat?
		\end{block}
		\begin{enumerate}[(A)]
			\footnotesize
			\item Background checks at hiring, repeated every few years, which screen out most malicious individuals
			\item A very strict network perimeter that blocks all unauthorized external access
			\item Least privilege, monitoring for anomalous behavior, segregation of duties, and HR-integrated onboarding and offboarding
			\item Prohibiting personal devices in the workplace so that data cannot be copied out of the organization
		\end{enumerate}
	\end{frame}
	
	\begin{frame}{Q397: Explanation}
		\begin{block}{Correct Answer: C}
			\textbf{Why (C) is the best answer:}
			\begin{itemize}
				\item Limit access to what is needed
				\item Monitor behavior
				\item Coordinate with HR at joiner, mover, leaver
			\end{itemize}
		\end{block}
		\begin{alertblock}{Critical Trap}
			Insider risk needs layered, people-centered controls.
		\end{alertblock}
	\end{frame}
	
	\begin{frame}{Q398: Limited Pen Test Scope}
		\begin{block}{Question}
			\small A penetration test excluded the most critical systems to avoid disruption. What is the MAIN implication?
		\end{block}
		\begin{enumerate}[(A)]
			\footnotesize
			\item The test is more reliable, because it avoided disruption and therefore produced cleaner results
			\item Critical systems are automatically secure, since they are subject to stricter change control
			\item The organization saved cost with no impact on risk, because the tested systems are representative of those excluded
			\item Assurance does not cover the highest-risk assets, so a controlled testing approach for them should be designed
		\end{enumerate}
	\end{frame}
	
	\begin{frame}{Q398: Explanation}
		\begin{block}{Correct Answer: D}
			\textbf{Why (D) is the best answer:}
			\begin{itemize}
				\item Scope determines assurance
				\item Critical assets need testing
				\item Use safe, controlled methods
			\end{itemize}
		\end{block}
		\begin{alertblock}{Critical Trap}
			Do not skip the assets that matter most.
		\end{alertblock}
	\end{frame}
	
	\begin{frame}{Q399: Access Management KRI}
		\begin{block}{Question}
			\small Which is the BEST key risk indicator for access management?
		\end{block}
		\begin{enumerate}[(A)]
			\footnotesize
			\item Number of password resets performed, which indicates the burden on the help desk
			\item Total number of users in the directory, which indicates the size of the attack surface
			\item Percentage of privileged accounts not recertified within policy, and number of orphaned accounts
			\item Number of help desk calls related to access, which indicates user satisfaction with access processes
		\end{enumerate}
	\end{frame}
	
	\begin{frame}{Q399: Explanation}
		\begin{block}{Correct Answer: C}
			\textbf{Why (C) is the best answer:}
			\begin{itemize}
				\item KRIs show risk trends
				\item Privileged accounts are high risk
				\item Orphaned accounts show process failure
			\end{itemize}
		\end{block}
		\begin{alertblock}{Critical Trap}
			A KRI measures risk exposure, not activity.
		\end{alertblock}
	\end{frame}
	
	\begin{frame}{Q400: Log Tampering}
		\begin{block}{Question}
			\small Administrators can edit the logs on the systems they manage. Which control is MOST appropriate?
		\end{block}
		\begin{enumerate}[(A)]
			\footnotesize
			\item Ask administrators to confirm in writing that they will not alter logs, and review the confirmation annually
			\item Forward logs in near real time to a centralized, access-restricted, tamper-evident repository
			\item Reduce log verbosity so that less information is available that can be edited
			\item Print logs weekly and file them for manual review by the internal audit function
		\end{enumerate}
	\end{frame}
	
	\begin{frame}{Q400: Explanation}
		\begin{block}{Correct Answer: B}
			\textbf{Why (B) is the best answer:}
			\begin{itemize}
				\item Logs must be independent of admins
				\item Centralize and restrict access
				\item Tamper evidence supports investigations
			\end{itemize}
		\end{block}
		\begin{alertblock}{Critical Trap}
			Logs controlled by those they record are not reliable evidence.
		\end{alertblock}
	\end{frame}
	
	\begin{frame}{Exam Mindset Summary}
		\begin{block}{Remember}
			\begin{itemize}
				\item \textbf{Judgment over trivia:} pick the action that best manages risk for the business
				\item \textbf{Governance first:} exceptions, deviations, and trade-offs need documented decisions by accountable owners
				\item \textbf{Know the dependency:} resilience, vendors, and shared services must match what depends on them
				\item \textbf{Test, do not assume:} restores, rollbacks, and compensating controls need evidence
				\item \textbf{Contain before you eradicate} and preserve evidence
				\item \textbf{Least privilege and independence:} no self-approval, no shared credentials
				\item \textbf{Data:} know where it is, who owns it, how long to keep it, and how to destroy it
				\item \textbf{Security is continuous risk management}
			\end{itemize}
		\end{block}
	\end{frame}
	
	\begin{frame}{Thank You}
		\begin{center}
			{\Huge \textbf{Good Luck on Your CRISC Exam!}}
		\end{center}
	\end{frame}
	
\end{document}
